\chapter{\texorpdfstring{\(\PP^2\) Reloaded}{P2 Reloaded}}
\label{chap:p2-reloaded}

\begin{chapteroverview}
In this chapter we revisit the projective plane from three perspectives:
its affine charts and homogeneous coordinate ring, its torus action,
and the combinatorics of a lattice triangle and its normal fan. We also see for
the first time that monomial sections of the line bundles
\(\cO_{\PP^2}(d)\) are indexed by the lattice points of a dilated
triangle. The goal is not yet to develop the general theory, but to
encounter its central correspondences in a familiar example.
\end{chapteroverview}

% Suggested learning goals:
% By the end of the chapter, students should be able to:
% 1. identify the dense torus in A^2 and P^2;
% 2. describe its characters and one-parameter subgroups;
% 3. recover the first-quadrant cone from extension properties;
% 4. describe the torus-orbit decomposition of P^2;
% 5. match degree-d monomials with lattice points in d Delta;
% 6. draw the fan of P^2;
% 7. recognize a toric blowup as a star subdivision.

\section{The affine model: \texorpdfstring{\(\AAA^2\)}{A2}}
\label{sec:p2-affine-model}
\index{affine plane}
\index{coordinate strata}

Our first example is the affine plane. Fix a field \(\Bbbk\), and write
\[
\AAA^2=\Spec \Bbbk[x,y].
\]
The symbols \(x\) and \(y\) are coordinate functions. Once they have
been chosen, every \(\Bbbk\)-point \(p\) determines a pair
\[
\bigl(x(p),y(p)\bigr)\in\Bbbk^2,
\]
and this gives the familiar identification
\[
\AAA^2(\Bbbk)\simeq\Bbbk^2.
\]
The notation \(\AAA^2\), however, remembers more than the set
\(\Bbbk^2\): it remembers which functions are polynomial and which
subsets are cut out by polynomial equations. This extra structure is
the algebraic geometry of the affine plane.

\begin{guidingquestion}
When \(\Bbbk=\CC\), what is the difference between
\(\AAA^2_{\CC}\) and \(\CC^2\)?
\end{guidingquestion}

Their sets of complex points are naturally identified. The notation
\(\CC^2\), however, often asks us to remember its structure as a
complex vector space or a complex manifold, together with the
Euclidean topology and holomorphic functions. The notation
\(\AAA^2_{\CC}\) asks us instead to remember the Zariski topology and
the polynomial functions. The points are the same; the chosen
geometric structure is different.

\subsection{The coordinate strata}

The coordinates \(x\) and \(y\) single out the two coordinate axes,
\[
V(x)=\{x=0\},
\qquad
V(y)=\{y=0\},
\]
and their intersection, the origin
\[
0=V(x,y).
\]
According to which coordinates vanish, the affine plane breaks into
four disjoint pieces:
\[
\AAA^2
=
T
\sqcup
\bigl(V(x)\setminus\{0\}\bigr)
\sqcup
\bigl(V(y)\setminus\{0\}\bigr)
\sqcup
\{0\},
\]
where
\[
T
=
\AAA^2\setminus V(xy)
=
\{(a,b):a\neq0,\ b\neq0\}.
\]

\begin{figure}[ht]
\centering
\begin{tikzpicture}[scale=0.95]
  \fill[DictionaryBlue!8] (-3,-2.3) rectangle (3,2.3);
  \draw[thick] (-3,0) -- (3,0);
  \draw[thick] (0,-2.3) -- (0,2.3);
  \fill (0,0) circle (1.6pt);

  \node at (1.55,1.25) {$T$};
  \node at (-1.55,1.25) {$T$};
  \node at (-1.55,-1.25) {$T$};
  \node at (1.55,-1.25) {$T$};

  \node[above] at (2.15,0) {$V(y)\setminus\{0\}$};
  \node[rotate=90,above] at (0,1.15) {$V(x)\setminus\{0\}$};
  \node[below left] at (0,0) {$0$};
\end{tikzpicture}
\caption{A schematic real-coordinate picture of the four strata of
\(\AAA^2\). The shading represents the open set \(T\); over a general
field, the picture should be read combinatorially rather than
literally.}
\label{fig:a2-strata}
\end{figure}

On \(T\), both coordinate functions are nonzero and hence invertible.
Its coordinate ring is therefore
\[
\Bbbk[T]=\Bbbk[x^{\pm1},y^{\pm1}],
\]
so
\[
T\simeq\Gm^2.
\]
We call \(T\) the \emph{dense torus} of \(\AAA^2\). It is open because
it is the complement of \(V(xy)\), and it is dense because
\(\AAA^2\) is irreducible and \(T\) is nonempty.

The name ``torus'' reflects the coordinatewise multiplication law
\[
(t_1,t_2)(s_1,s_2)=(t_1s_1,t_2s_2).
\]
Using this law, \(T\) acts on the whole affine plane by rescaling:
\[
(t_1,t_2)\cdot(a,b)=(t_1a,t_2b).
\]
The four pieces above are exactly the four orbits of this action.
The open set \(T\) is the largest orbit; each punctured coordinate
axis is a one-dimensional orbit; and the origin is fixed. Thus the
boundary of the torus already comes naturally organized into strata.

\subsection{Monomial functions and lattice points}

The polynomial functions most naturally adapted to coordinatewise
rescaling are the monomials
\[
x^a y^b,
\qquad
a,b\in\ZZ_{\geq0}.
\]
Every such monomial is regular on \(\AAA^2\), and none vanishes at a
point of \(T\). Its two exponents provide a compact combinatorial
label, so we record the monomial \(x^ay^b\) by the lattice point
\[
(a,b)\in M:=\ZZ^2.
\]

The monomials in \(\Bbbk[x,y]\) occupy precisely the lattice points in
the nonnegative quadrant:
\[
\ZZ_{\geq0}^2\subseteq M_{\RR}\simeq\RR^2.
\]
The corresponding real cone is
\[
\sigma^\vee=\RR_{\geq0}^2.
\]

\begin{figure}[ht]
\centering
\begin{tikzpicture}[scale=0.9]
  \fill[SurfaceGreen!10] (0,0) rectangle (3.5,3.5);
  \draw[->] (-0.2,0) -- (3.9,0) node[right] {$a$};
  \draw[->] (0,-0.2) -- (0,3.9) node[above] {$b$};

  \foreach \i in {0,1,2,3}
    \foreach \j in {0,1,2,3}
      \fill (\i,\j) circle (1.2pt);

  \draw[ray] (0,0) -- (3.6,0);
  \draw[ray] (0,0) -- (0,3.6);

  \node at (2.65,3.0) {$\sigma^\vee$};
  \node[below right] at (0,0) {$1$};
  \node[below] at (1,0) {$x$};
  \node[left] at (0,1) {$y$};
  \node[above right] at (2,1) {$x^2y$};
  \node[above right] at (1,2) {$xy^2$};
\end{tikzpicture}
\caption{The lattice point \((a,b)\) labels the monomial \(x^ay^b\).
The lattice points of \(\sigma^\vee\) are exactly the monomials regular
on \(\AAA^2\).}
\label{fig:a2-monomials-cone}
\end{figure}

The significance of the quadrant becomes clearer if we begin on the
dense torus. Since \(x\) and \(y\) are invertible on \(T\), every
Laurent monomial
\[
x^a y^b,
\qquad
(a,b)\in\ZZ^2,
\]
is regular there. Such a Laurent monomial extends across the boundary
to a regular function on all of \(\AAA^2\) precisely when
\[
a\geq0
\qquad\text{and}\qquad
b\geq0.
\]
Indeed, a negative exponent of \(x\) produces a pole along \(V(x)\),
and a negative exponent of \(y\) produces a pole along \(V(y)\).

This gives a useful way to recover the affine plane from its dense
torus: retain exactly those Laurent monomials whose exponent vectors
lie in \(\sigma^\vee\). In symbols,
\[
\Bbbk[x,y]
=
\bigoplus_{(a,b)\in\sigma^\vee\cap M}
\Bbbk\,x^ay^b.
\]
The cone is therefore not merely a picture attached to the coordinate
ring; it records exactly which monomial functions survive when the
boundary is added.

\begin{observation}
In the representation theory of algebraic tori, a Laurent monomial
\(x^ay^b\) is called a \emph{character} of \(T\), and \(M\) is called
the \emph{character lattice}. No representation theory is required
for the discussion here. For now, \(M\) may simply be regarded as the
lattice indexing Laurent monomials on \(T\).
\end{observation}

\subsection{Approaching the boundary along monomial curves}

There is another first quadrant, obtained by looking at simple curves
in the torus rather than functions on it. For an integer pair
\((u,v)\in\ZZ^2\), consider
\[
\Gm\longrightarrow T,
\qquad
t\longmapsto(t^u,t^v).
\]
We ask whether this curve acquires a value at \(t=0\). Algebraically,
this means asking whether the map extends to a morphism
\[
\AAA^1\longrightarrow\AAA^2.
\]

The extension exists exactly when
\[
u\geq0
\qquad\text{and}\qquad
v\geq0.
\]
If both exponents are nonnegative, the two coordinates \(t^u\) and
\(t^v\) are polynomial functions of \(t\). If either exponent is
negative, the corresponding coordinate has a pole at \(t=0\).

The position of \((u,v)\) in the quadrant determines where the
extended curve meets the boundary:
\[
\begin{array}{c|c}
\text{position of }(u,v) & \text{value at }t=0\\
\hline
u>0,\ v>0 & (0,0)\\
u>0,\ v=0 & (0,1)\in V(x)\setminus\{0\}\\
u=0,\ v>0 & (1,0)\in V(y)\setminus\{0\}\\
u=v=0 & (1,1)\in T.
\end{array}
\]
Thus the relative interior of the two-dimensional quadrant points
toward the origin, the relative interior of each ray points toward
one of the punctured coordinate axes, and the zero-dimensional face
points toward the torus itself.

The particular points \((0,1)\), \((1,0)\), and \((1,1)\) arise
because our curve begins at the identity of the torus. More generally,
for \((c_1,c_2)\in T\), consider the translated curve
\[
t\longmapsto(c_1t^u,c_2t^v).
\]
If \(u>0\) and \(v=0\), its value at \(0\) is \((0,c_2)\), and every
point of \(V(x)\setminus\{0\}\) occurs in this way. Similarly, the
other ray produces all of \(V(y)\setminus\{0\}\). The interior always
produces the origin. In this precise sense, the faces of the quadrant
encode the four strata of \(\AAA^2\).

\begin{figure}[ht]
\centering
\begin{tikzpicture}[scale=1]
  \fill[ExampleGold!10] (0,0) rectangle (3.7,3.7);
  \draw[->] (-0.2,0) -- (4.1,0) node[right] {$u$};
  \draw[->] (0,-0.2) -- (0,4.1) node[above] {$v$};
  \draw[ray] (0,0) -- (3.8,0);
  \draw[ray] (0,0) -- (0,3.8);

  \draw[thick,->] (0,0) -- (2.7,1.9);
  \node[above right,align=left] at (2.55,2.0)
    {interior\\[-1mm]\(\leadsto\{0\}\)};

  \node[below,align=center] at (2.15,0)
    {ray\\[-1mm]\(\leadsto V(x)\setminus\{0\}\)};
  \node[rotate=90,above,align=center] at (0,1.85)
    {ray \(\leadsto V(y)\setminus\{0\}\)};

  \fill (0,0) circle (1.5pt);
  \node[below left,align=right] at (0,0)
    {zero face\\[-1mm]\(\leadsto T\)};
\end{tikzpicture}
\caption{Faces of the cone of extendable monomial curves correspond to
the torus-orbit strata of \(\AAA^2\). The arrows indicate the exponent
vectors \((u,v)\), not paths in the affine plane.}
\label{fig:a2-limit-cone}
\end{figure}

The exponent vectors of the extendable curves form the cone
\[
\sigma=\RR_{\geq0}^2\subseteq N_{\RR},
\qquad
N:=\ZZ^2.
\]
Although \(\sigma\) and \(\sigma^\vee\) are both drawn as first
quadrants in our chosen coordinates, they answer different questions:
\(\sigma^\vee\) records monomial functions that extend to
\(\AAA^2\), while \(\sigma\) records monomial curves that extend to
\(\AAA^2\).

The connection between the two is visible in one calculation. Evaluate
the monomial \(x^ay^b\) along the curve \(t\mapsto(t^u,t^v)\):
\[
x^ay^b(t^u,t^v)=t^{au+bv}.
\]
The integer \(au+bv\) records the behavior of the function along the
curve. A positive value means that it vanishes at \(t=0\), a zero
value means that it remains nonzero, and a negative value means that
it has a pole.

\begin{observation}
The curves \(t\mapsto(t^u,t^v)\) are called
\emph{one-parameter subgroups} of \(T\), and \(N\) is called the
\emph{lattice of one-parameter subgroups}. The integer \(au+bv\) is
written
\[
\pair{(a,b)}{(u,v)}.
\]
This is the natural pairing between the two exponent lattices \(M\)
and \(N\). The terminology is useful, but the essential content is
the elementary identity
\[
x^ay^b(t^u,t^v)=t^{au+bv}.
\]
\end{observation}

We can now state the duality precisely:
\[
\sigma^\vee
=
\left\{
m\in M_{\RR}:
\pair{m}{n}\geq0
\text{ for every }n\in\sigma
\right\}.
\]
Indeed, a monomial indexed by \(m\) extends across every monomial curve
indexed by \(\sigma\) exactly when all of these pairings are
nonnegative. The two cones look alike in the chosen coordinates, but
they live in paired vector spaces and encode complementary aspects of
the affine plane.

\begin{threedictionaries}
For the affine plane, our first dictionary is
\begin{center}
\small
\begin{tabularx}{\linewidth}{
  >{\centering\arraybackslash}X|
  >{\centering\arraybackslash}X|
  >{\centering\arraybackslash}X}
\textbf{Algebra} & \textbf{Geometry} & \textbf{Combinatorics}\\
\hline
\(x^ay^b\in\Bbbk[x,y]\) &
a monomial extending across the boundary &
\((a,b)\in\sigma^\vee\cap M\)
\\[2mm]
\(t\mapsto(t^u,t^v)\) &
a curve acquiring a boundary value &
\((u,v)\in\sigma\cap N\)
\\[2mm]
\(x^ay^b(t^u,t^v)=t^{au+bv}\) &
order along the curve &
\(\pair{(a,b)}{(u,v)}=au+bv\).
\end{tabularx}
\end{center}
\end{threedictionaries}


\section{What we already know about
\texorpdfstring{\(\PP^2\)}{P2}}
\label{sec:p2-review}
\index{projective plane}
\index{homogeneous coordinates}
\index{coordinate triangle}

We now pass from the affine plane to the projective plane. The point of
this review is not merely to repeat definitions we already know.
Instead, we isolate a few familiar features of \(\PP^2\) that will
soon acquire a toric interpretation: its three affine charts, its
homogeneous coordinates, and the triangle formed by the coordinate
lines.

\subsection{Lines through the origin}

For the moment, think of the \(\Bbbk\)-points of \(\PP^2\) as
one-dimensional linear subspaces of \(\Bbbk^3\). A nonzero vector
\[
(x_0,x_1,x_2)\in\Bbbk^3
\]
determines the line that it spans, and multiplying the vector by a
nonzero scalar does not change that line. Thus
\[
\PP^2(\Bbbk)
=
\frac{\Bbbk^3\setminus\{0\}}{\Bbbk^\times}.
\]
We write the equivalence class of \((x_0,x_1,x_2)\) as
\[
[x_0:x_1:x_2].
\]
These are the \emph{homogeneous coordinates} on \(\PP^2\). By
definition,
\[
[x_0:x_1:x_2]
=
[\lambda x_0:\lambda x_1:\lambda x_2]
\qquad
\text{for every }\lambda\in\Bbbk^\times,
\]
and the triple \((0,0,0)\) is not permitted.

A projective point therefore records a direction in \(\Bbbk^3\), rather
than a vector with a specified size. One may choose a convenient
representative of the direction, but there is usually no preferred
representative valid for every point at once.

\begin{figure}[ht]
\centering
\begin{tikzpicture}[x=1cm,y=1cm]
  % In this two-dimensional schematic, the horizontal coordinate
  % represents x_0 and all remaining directions are compressed into
  % the vertical coordinate.
  \coordinate (O) at (0,0);

  \fill[PastelYellow] (4,-2.0) rectangle (4.35,2.0);
  \draw[thick,ExampleGold] (4,-2.0) -- (4,2.0);
  \node[above,ExampleGold!70!black] at (4.15,2.0)
    {\(\Pi_0=\{x_0=1\}\)};

  \draw[thick,DictionaryBlue] (O) -- (4,-1.25);
  \draw[thick,SurfaceGreen] (O) -- (4,0.15);
  \draw[thick,ExampleGold] (O) -- (4,1.35);

  \fill[DictionaryBlue] (4,-1.25) circle (1.8pt);
  \fill[SurfaceGreen] (4,0.15) circle (1.8pt);
  \fill[ExampleGold] (4,1.35) circle (1.8pt);

  \draw[very thick,WarningRed] (0,-1.75) -- (0,1.75);
  \node[left,align=right,WarningRed] at (-0.15,1.25)
    {directions with\\\(x_0=0\)};

  \fill (O) circle (1.8pt);
  \node[below left] at (O) {\(0\)};

  \draw[->,black!55] (0,-2.25) -- (4.8,-2.25)
    node[right] {\(x_0\)};
\end{tikzpicture}
\caption{Projectivization in a two-dimensional schematic. A line
through the origin with \(x_0\neq0\) meets the normalization slice
\(\Pi_0=\{x_0=1\}\) exactly once. A direction contained in
\(\{x_0=0\}\) never meets that slice.}
\label{fig:p2-lines-through-origin}
\end{figure}

\subsection{The standard affine charts}

Suppose first that \(x_0\neq0\). We may rescale the homogeneous
coordinates by \(x_0^{-1}\) and write the point uniquely in the form
\[
[1:z_1:z_2],
\qquad
z_1=\frac{x_1}{x_0},
\quad
z_2=\frac{x_2}{x_0}.
\]
This identifies the open subset
\[
U_0=\{[x_0:x_1:x_2]\in\PP^2:x_0\neq0\}
\]
with the affine plane:
\[
\begin{aligned}
\phi_0:U_0&\longrightarrow\AAA^2,\\
[x_0:x_1:x_2]
&\longmapsto
\left(\frac{x_1}{x_0},\frac{x_2}{x_0}\right).
\end{aligned}
\]
The inverse map is
\[
(a,b)\longmapsto[1:a:b].
\]

There is nothing special about \(x_0\). Repeating the construction
with the other two coordinates gives
\[
U_i=\{x_i\neq0\}\simeq\AAA^2,
\qquad
i=0,1,2.
\]
Every projective point has at least one nonzero coordinate, so
\[
\PP^2=U_0\cup U_1\cup U_2.
\]
Thus \(\PP^2\) is covered by three copies of the affine plane.

\begin{figure}[ht]
\centering
\begin{tikzpicture}[scale=0.98]
  \coordinate (p0) at (-1.7,3.05);
  \coordinate (p1) at (-4.35,0);
  \coordinate (p2) at (0.95,0);

  \fill[PastelYellow] (p0) -- (p1) -- (p2) -- cycle;

  \draw[very thick,SurfaceGreen] (p0) -- (p1);
  \draw[very thick,DictionaryBlue] (p0) -- (p2);
  \draw[very thick,WarningRed] (p1) -- (p2);

  \fill[ExampleGold] (p0) circle (2.2pt);
  \fill[DictionaryBlue] (p1) circle (2.2pt);
  \fill[SurfaceGreen] (p2) circle (2.2pt);

  \node[above] at (p0) {\(p_0=[1:0:0]\)};
  \node[below left] at (p1) {\(p_1=[0:1:0]\)};
  \node[below right] at (p2) {\(p_2=[0:0:1]\)};

  \node[fill=white,rounded corners,inner sep=4pt,align=center]
    at (-1.7,1.35)
    {\(U_0=\PP^2\setminus D_0\)\\[-1mm]
     \(\simeq\AAA^2\)};

  % Legend placed to the right, away from the triangle.
  \node[anchor=west,font=\small] at (2.0,2.35) {coordinate lines};

  \draw[very thick,WarningRed] (2.0,1.75) -- (2.65,1.75);
  \node[anchor=west] at (2.82,1.75) {\(D_0=\{x_0=0\}\)};

  \draw[very thick,DictionaryBlue] (2.0,1.05) -- (2.65,1.05);
  \node[anchor=west] at (2.82,1.05) {\(D_1=\{x_1=0\}\)};

  \draw[very thick,SurfaceGreen] (2.0,0.35) -- (2.65,0.35);
  \node[anchor=west] at (2.82,0.35) {\(D_2=\{x_2=0\}\)};
\end{tikzpicture}
\caption{The coordinate triangle. The three colors distinguish the
coordinate lines. In this schematic, the pale yellow region indicates
the chart \(U_0\), obtained by removing the red side \(D_0\).}
\label{fig:p2-coordinate-triangle-preview}
\end{figure}

The complement of \(U_0\) is the projective line
\[
D_0=\{x_0=0\}\simeq\PP^1.
\]
Relative to the chart \(U_0\), this is the \emph{line at infinity}.
Hence one may think schematically of
\[
\PP^2
=
\AAA^2\ \sqcup\ \PP^1,
\]
where the affine plane is \(U_0\) and the projective line is \(D_0\).
This description depends on the choice of \(x_0\): choosing another
chart produces another affine plane and another line at infinity.

\subsection{Changing affine coordinates}

The affine charts overlap, and their coordinates are related by
explicit formulas. On \(U_0\), write
\[
[x_0:x_1:x_2]=[1:z_1:z_2].
\]
On \(U_1\), write the same point as
\[
[x_0:x_1:x_2]=[w_0:1:w_2].
\]
A point lies in \(U_0\cap U_1\) exactly when \(z_1\neq0\). To pass
from the \(U_0\)-coordinates to the \(U_1\)-coordinates, we simply
renormalize so that the second homogeneous coordinate becomes \(1\):
\[
[1:z_1:z_2]
=
\left[\frac1{z_1}:1:\frac{z_2}{z_1}\right].
\]
Thus the transition map is
\[
(z_1,z_2)
\longmapsto
\left(
\frac1{z_1},
\frac{z_2}{z_1}
\right).
\]
The inverse is obtained by the same kind of renormalization:
\[
(w_0,w_2)
\longmapsto
\left(
\frac1{w_0},
\frac{w_2}{w_0}
\right).
\]

This simple calculation contains an important lesson about projective
space. Each chart has two ordinary
coordinates, but no single pair of affine coordinates works
everywhere. On overlaps, the coordinates are related by rational
expressions whose denominators are nonzero precisely where the two
charts meet.

\begin{figure}[ht]
\centering
\begin{tikzpicture}[
  source/.style={
    rounded corners,
    draw=black!60,
    fill=black!3,
    minimum width=3.4cm,
    minimum height=1.0cm,
    align=center
  },
  chartzero/.style={
    rounded corners,
    draw=ExampleGold!80!black,
    fill=PastelYellow,
    minimum width=3.35cm,
    minimum height=1.45cm,
    align=center
  },
  chartone/.style={
    rounded corners,
    draw=DictionaryBlue,
    fill=DictionaryBlue!7,
    minimum width=3.35cm,
    minimum height=1.45cm,
    align=center
  },
  charttwo/.style={
    rounded corners,
    draw=SurfaceGreen,
    fill=SurfaceGreen!8,
    minimum width=3.35cm,
    minimum height=1.45cm,
    align=center
  }
]
  \node[source] (p) at (0,3.0)
    {\(\displaystyle [x_0:x_1:x_2]\)};

  \node[chartzero] (U0) at (-4.0,0)
    {\textcolor{ExampleGold!80!black}{\(x_0\neq0\): divide by \(x_0\)}\\[1mm]
     \(\displaystyle \left[1:\frac{x_1}{x_0}:\frac{x_2}{x_0}\right]\)};

  \node[chartone] (U1) at (0,0)
    {\textcolor{DictionaryBlue}{\(x_1\neq0\): divide by \(x_1\)}\\[1mm]
     \(\displaystyle \left[\frac{x_0}{x_1}:1:\frac{x_2}{x_1}\right]\)};

  \node[charttwo] (U2) at (4.0,0)
    {\textcolor{SurfaceGreen}{\(x_2\neq0\): divide by \(x_2\)}\\[1mm]
     \(\displaystyle \left[\frac{x_0}{x_2}:\frac{x_1}{x_2}:1\right]\)};

  \draw[thick,-{Latex[length=2.2mm]},ExampleGold!80!black]
    (p) -- (U0);
  \draw[thick,-{Latex[length=2.2mm]},DictionaryBlue]
    (p) -- (U1);
  \draw[thick,-{Latex[length=2.2mm]},SurfaceGreen]
    (p) -- (U2);
\end{tikzpicture}
\caption{The three affine charts arise from the same operation:
choose a nonzero homogeneous coordinate and normalize it to \(1\).
Transition maps are obtained by normalizing the same point in a
different way.}
\label{fig:p2-chart-transitions}
\end{figure}

\subsection{Homogeneous equations}

Because homogeneous coordinates are defined only up to a common
nonzero scalar, not every polynomial expression gives a well-defined
condition on projective points. The appropriate equations are
homogeneous polynomials.

A polynomial \(F(x_0,x_1,x_2)\) is homogeneous of degree \(d\) if
\[
F(\lambda x_0,\lambda x_1,\lambda x_2)
=
\lambda^dF(x_0,x_1,x_2)
\]
for every scalar \(\lambda\). It follows that
\[
F(x_0,x_1,x_2)=0
\]
is independent of the chosen representative of
\([x_0:x_1:x_2]\). For example, a homogeneous linear equation cuts
out a projective line, while a homogeneous quadratic equation cuts
out a projective conic.

The polynomial ring
\[
S=\Bbbk[x_0,x_1,x_2]
\]
is naturally graded by total degree:
\[
S=\bigoplus_{d\geq0}S_d,
\]
where \(S_d\) is the vector space of homogeneous polynomials of degree
\(d\). This is the \emph{homogeneous coordinate ring} of \(\PP^2\).
The grading is essential: it records how an expression changes when
the homogeneous coordinates are simultaneously rescaled.

\begin{observation}
The phrase ``coordinate ring'' is being used differently here than it
was for \(\AAA^2\). The ring \(\Bbbk[x,y]\) is the ring of global
regular functions on \(\AAA^2\). By contrast,
\(\Bbbk[x_0,x_1,x_2]\) is not the ring of global regular functions on
\(\PP^2\); its grading encodes projective equations and, later,
sections of line bundles. This distinction is one reason homogeneous
coordinates are richer than ordinary affine coordinates.
\end{observation}

\subsection{The coordinate triangle}

The three homogeneous coordinate functions define three projective
lines
\[
D_0=V(x_0),
\qquad
D_1=V(x_1),
\qquad
D_2=V(x_2).
\]
We call their union
\[
D_0\cup D_1\cup D_2
=
V(x_0x_1x_2)
\]
the \emph{coordinate triangle}. Its vertices are the three coordinate
points
\[
p_0=[1:0:0],
\qquad
p_1=[0:1:0],
\qquad
p_2=[0:0:1].
\]
More precisely,
\[
p_0=D_1\cap D_2,
\qquad
p_1=D_0\cap D_2,
\qquad
p_2=D_0\cap D_1.
\]
There is no common intersection of all three coordinate lines, because
\([0:0:0]\) is not a projective point.

The coordinate triangle organizes \(\PP^2\) according to which
homogeneous coordinates vanish:
\[
\begin{array}{c|c}
\text{vanishing pattern} & \text{resulting piece}\\
\hline
x_0x_1x_2\neq0
  & \text{one open piece}\\
\text{exactly one }x_i=0
  & \text{three punctured coordinate lines}\\
\text{exactly two }x_i=0
  & \text{the three coordinate points}.
\end{array}
\]
Thus the coordinate decomposition of \(\PP^2\) has seven pieces. In
the next section, we will see that these are not merely convenient
subsets: they are the natural strata for a two-dimensional torus
action.

\begin{guidingquestion}
What remains after the coordinate triangle
\(V(x_0x_1x_2)\) is removed from \(\PP^2\), and how is it related to
the dense torus in \(\AAA^2\)?
\end{guidingquestion}


\section{The dense torus in
\texorpdfstring{\(\PP^2\)}{P2}}
\label{sec:p2-dense-torus}
\index{torus!dense torus}
\index{torus action!on the projective plane}
\index{projective plane!dense torus}

At the end of the previous section, we decomposed \(\PP^2\) according
to which homogeneous coordinates vanish. The largest piece is the one
on which no coordinate vanishes:
\[
T
=
\PP^2\setminus V(x_0x_1x_2)
=
\bigl\{[x_0:x_1:x_2]:x_0x_1x_2\neq0\bigr\}.
\]
We now take a closer look at this open set. It will play the same role
in \(\PP^2\) that \(\Gm^2\) played inside \(\AAA^2\).

Throughout this chapter, the phrase \emph{coordinate triangle} means
the union
\[
D_0\cup D_1\cup D_2
=
V(x_0x_1x_2)
\]
of the three coordinate lines. We are not asserting that this union is
a triangle in any Euclidean metric on \(\PP^2\). The terminology is
combinatorial: the three lines meet pairwise in three points, producing
the incidence pattern of three sides and three vertices.

\subsection{Removing the coordinate triangle}

The condition \(x_0x_1x_2\neq0\) means that all three homogeneous
coordinates are nonzero. In particular, every point of \(T\) lies in
the chart \(U_0\). After dividing by \(x_0\), it has a unique
representative
\[
[1:z_1:z_2]
\qquad\text{with}\qquad
z_1,z_2\in\Bbbk^\times.
\]
Therefore the map
\[
\begin{aligned}
T&\longrightarrow\Gm^2,\\
[x_0:x_1:x_2]
&\longmapsto
\left(\frac{x_1}{x_0},\frac{x_2}{x_0}\right)
\end{aligned}
\]
is an isomorphism, with inverse
\[
(z_1,z_2)\longmapsto[1:z_1:z_2].
\]
Thus
\[
T\simeq\Gm^2.
\]

\begin{figure}[ht]
\centering
\begin{tikzpicture}[scale=1.0]
  \coordinate (p0) at (0,3.35);
  \coordinate (p1) at (-3.15,0);
  \coordinate (p2) at (3.15,0);

  \fill[DictionaryBlue!10] (p0) -- (p1) -- (p2) -- cycle;

  % Colored boundary.
  \draw[very thick,SurfaceGreen] (p0) -- (p1);
  \draw[very thick,ExampleGold] (p0) -- (p2);
  \draw[very thick,WarningRed] (p1) -- (p2);

  % Open torus region.
  \fill[PastelYellow,opacity=0.95]
    (-0.15,2.55) -- (-2.15,0.42) -- (2.15,0.42) -- (0.15,2.55) -- cycle;
  \node[align=center] at (0,1.45)
    {\(\displaystyle T=\{x_0x_1x_2\neq0\}\)\\[1mm]
     \(\displaystyle \simeq\Gm^2\)};

  \fill[DictionaryBlue] (p0) circle (2pt);
  \fill[SurfaceGreen] (p1) circle (2pt);
  \fill[ExampleGold] (p2) circle (2pt);

  \node[above] at (p0) {\(p_0\)};
  \node[below left] at (p1) {\(p_1\)};
  \node[below right] at (p2) {\(p_2\)};

  \begin{scope}[yshift=-0.85cm]
    \draw[very thick,WarningRed] (-3.0,0) -- (-2.45,0);
    \node[right] at (-2.38,0) {\(D_0\)};

    \draw[very thick,ExampleGold] (-0.55,0) -- (0,0);
    \node[right] at (0.07,0) {\(D_1\)};

    \draw[very thick,SurfaceGreen] (1.75,0) -- (2.30,0);
    \node[right] at (2.37,0) {\(D_2\)};
  \end{scope}
\end{tikzpicture}
\caption{The dense torus is what remains after the three coordinate
lines are removed. The drawing is only a combinatorial model: the
``triangle'' means the union \(D_0\cup D_1\cup D_2\), not a Euclidean
triangle sitting inside \(\PP^2\). The yellow region represents the
single open set on which all three coordinates are nonzero.}
\label{fig:p2-dense-torus-region}
\end{figure}

This description also explains why the word \emph{dense} is
appropriate. The complement of \(T\) is the union of the three
coordinate lines---what we are calling the coordinate triangle---and
this union is a proper closed subset of \(\PP^2\). Since \(\PP^2\) is
irreducible, every nonempty open subset is dense.

\begin{guidingquestion}
Why do three nonzero homogeneous coordinates produce only a
two-dimensional torus?
\end{guidingquestion}

The answer is that homogeneous coordinates contain one redundant
degree of freedom. A triple
\[
(x_0,x_1,x_2)\in(\Bbbk^\times)^3
\]
and its scalar multiple
\[
(\lambda x_0,\lambda x_1,\lambda x_2)
\]
represent the same projective point. Dividing by the common scalar
leaves two independent ratios. On \(U_0\), those ratios are
\[
z_1=\frac{x_1}{x_0},
\qquad
z_2=\frac{x_2}{x_0}.
\]

\begin{figure}[ht]
\centering
\begin{tikzpicture}[
  box/.style={
    rounded corners,
    minimum height=1.25cm,
    minimum width=3.2cm,
    align=center,
    draw=black!60
  }
]
  \node[box,fill=DictionaryBlue!8] (left) at (-3.8,0)
    {\(\displaystyle [x_0:x_1:x_2]\in T\)\\
     \(x_0x_1x_2\neq0\)};

  \node[box,fill=PastelYellow] (middle) at (0,0)
    {\(\displaystyle
      [1:x_1/x_0:x_2/x_0]\)\\
     normalize \(x_0\) to \(1\)};

  \node[box,fill=SurfaceGreen!9] (right) at (3.8,0)
    {\(\displaystyle
      (z_1,z_2)\in\Gm^2\)};

  \draw[thick,-{Latex[length=2.2mm]},DictionaryBlue]
    (left) -- (middle);

  \draw[thick,-{Latex[length=2.2mm]},SurfaceGreen]
    (middle) -- (right);

  \draw[thick,-{Latex[length=2.2mm]},ExampleGold]
    (right.south) to[out=-110,in=-70]
    node[below=2mm,fill=white,inner sep=2pt]
    {\(\displaystyle (z_1,z_2)\longmapsto[1:z_1:z_2]\)}
    (left.south);
\end{tikzpicture}
\caption{The identification \(T\simeq\Gm^2\). Moving from left to
right means normalizing \(x_0\) to \(1\) and retaining the two
independent ratios.}
\label{fig:p2-torus-identification}
\end{figure}

\subsection{Multiplication inside the open set}

The identification \(T\simeq\Gm^2\) transports the usual
coordinatewise multiplication on \(\Gm^2\) to \(T\). In normalized
coordinates,
\[
[1:z_1:z_2]\cdot[1:w_1:w_2]
=
[1:z_1w_1:z_2w_2].
\]
The identity element is
\[
e=[1:1:1],
\]
and
\[
[1:z_1:z_2]^{-1}
=
[1:z_1^{-1}:z_2^{-1}].
\]
So the open subset \(T\) is not merely isomorphic to a torus as a
variety: it carries a natural algebraic group structure.

There is also a symmetric way to write this multiplication without
choosing \(x_0\) as the preferred coordinate. If
\[
p=[x_0:x_1:x_2],
\qquad
q=[y_0:y_1:y_2]
\]
belong to \(T\), define
\[
p\cdot q
=
[x_0y_0:x_1y_1:x_2y_2].
\]
This is well defined. Indeed, rescaling \(p\) or \(q\) only rescales
all three products by a common nonzero scalar. The identity for this
description is again \([1:1:1]\).

\begin{observation}
A more representation-theoretic description is
\[
T\simeq\Gm^3/\Gm,
\]
where the copy of \(\Gm\) in the denominator is embedded diagonally:
\[
\lambda\longmapsto(\lambda,\lambda,\lambda).
\]
This says exactly that three nonzero coordinates are being considered
up to simultaneous rescaling. Nothing in this chapter requires the
language of quotient groups; the two affine ratios
\(x_1/x_0\) and \(x_2/x_0\) contain the same information.
\end{observation}

\subsection{Letting the torus move the whole projective plane}

The multiplication on \(T\) extends to an action on all of
\(\PP^2\). Using the coordinates \(T\simeq\Gm^2\), define
\[
(t_1,t_2)\cdot[x_0:x_1:x_2]
=
[x_0:t_1x_1:t_2x_2].
\]
This formula makes sense even when some of the \(x_i\) vanish. It is
also compatible with projective scaling: replacing
\([x_0:x_1:x_2]\) by
\([\lambda x_0:\lambda x_1:\lambda x_2]\) changes the output by the
same common scalar \(\lambda\).

On the open set \(T\), this action is simply multiplication. Indeed,
\[
(t_1,t_2)\cdot[1:z_1:z_2]
=
[1:t_1z_1:t_2z_2].
\]
In particular, every point of \(T\) can be moved to every other point
of \(T\). Given
\[
[1:z_1:z_2],
\qquad
[1:w_1:w_2]\in T,
\]
the unique torus element carrying the first point to the second is
\[
\left(\frac{w_1}{z_1},\frac{w_2}{z_2}\right).
\]
Thus \(T\) is a single orbit for the action.

\begin{figure}[ht]
\centering
\begin{tikzpicture}[
  pointbox/.style={
    rounded corners,
    draw=black!55,
    minimum width=3.35cm,
    minimum height=1.15cm,
    align=center
  }
]
  \node[pointbox,fill=DictionaryBlue!8] (p) at (-3.9,0)
    {\(\displaystyle p=[1:z_1:z_2]\)};

  \node[pointbox,fill=PastelYellow] (t) at (0,1.85)
    {\(\displaystyle t=(t_1,t_2)\in T\)};

  \node[pointbox,fill=SurfaceGreen!9] (q) at (3.9,0)
    {\(\displaystyle t\cdot p
      =[1:t_1z_1:t_2z_2]\)};

  \draw[very thick,-{Latex[length=2.4mm]},ExampleGold]
    (p) -- node[above=2mm,fill=white,inner sep=2pt]
    {\(\displaystyle \cdot\,(t_1,t_2)\)}
    (q);

  \draw[thick,-{Latex[length=2mm]},DictionaryBlue]
    (t) -- (0,0.62);
\end{tikzpicture}
\caption{The torus action on its open orbit is coordinatewise
multiplication. The same formula continues to make sense on the
coordinate lines and coordinate points.}
\label{fig:p2-torus-action}
\end{figure}

The action preserves each coordinate line. For example, if \(x_1=0\),
then
\[
(t_1,t_2)\cdot[x_0:0:x_2]
=
[x_0:0:t_2x_2],
\]
so the second coordinate remains zero. The three coordinate points
are fixed:
\[
(t_1,t_2)\cdot p_i=p_i
\qquad
(i=0,1,2).
\]
The complete orbit decomposition appears in
\cref{sec:p2-orbits}. For now, the essential picture is that the torus
moves freely through the large open set, while the coordinate triangle
records the ways in which the orbit can degenerate.

\subsection{The affine torus reappears in every chart}

Although we used \(U_0\) to identify \(T\) with \(\Gm^2\), the same
open torus appears in every standard affine chart. In \(U_1\), it is
described by
\[
\frac{x_0}{x_1}\neq0,
\qquad
\frac{x_2}{x_1}\neq0,
\]
and in \(U_2\) by
\[
\frac{x_0}{x_2}\neq0,
\qquad
\frac{x_1}{x_2}\neq0.
\]
Thus each chart contains the same open torus, but uses a different
pair of coordinates on it.

On \(U_0\cap T\), write
\[
z_1=\frac{x_1}{x_0},
\qquad
z_2=\frac{x_2}{x_0}.
\]
On \(U_1\cap T\), write
\[
w_0=\frac{x_0}{x_1},
\qquad
w_2=\frac{x_2}{x_1}.
\]
Then
\[
w_0=z_1^{-1},
\qquad
w_2=z_1^{-1}z_2.
\]
Notice that the change of coordinates is expressed entirely using
Laurent monomials. This is our first hint that monomials are not an
artifact of one chosen affine chart: they are the natural coordinate
changes on the torus itself.

\begin{threedictionaries}
The dense open set of \(\PP^2\) may be read in three parallel ways.
\begin{center}
\small
\begin{tabularx}{\linewidth}{
  >{\centering\arraybackslash}X|
  >{\centering\arraybackslash}X|
  >{\centering\arraybackslash}X}
\textbf{Algebra} & \textbf{Geometry} & \textbf{Combinatorics}\\
\hline
\(\Bbbk[z_1^{\pm1},z_2^{\pm1}]\) &
\(T=\PP^2\setminus V(x_0x_1x_2)\) &
two independent exponent directions
\\[2mm]
Laurent monomial coordinate changes &
the same torus inside all three affine charts &
integral changes of lattice coordinates
\\[2mm]
coordinatewise multiplication &
a torus action on \(\PP^2\) &
rescaling in two independent directions.
\end{tabularx}
\end{center}
\end{threedictionaries}

\begin{guidingquestion}
The action fixes the three coordinate points and preserves each
coordinate line. What are the remaining orbits on the coordinate
triangle?
\end{guidingquestion}


\section{Monomials, triangles, and global sections}
\label{sec:p2-characters-sections}
\index{monomial}
\index{global section}
\index{lattice point!and monomial sections}
\index{polytope!triangle of monomials}

We have now met two kinds of coordinate-like objects on \(\PP^2\).
Homogeneous coordinates
\[
[x_0:x_1:x_2]
\]
treat the three coordinate directions symmetrically, while affine
coordinates such as
\[
z_1=\frac{x_1}{x_0},
\qquad
z_2=\frac{x_2}{x_0}
\]
describe one chart at a time.

The qualifier ``coordinate-like'' matters. The symbols
\([x_0:x_1:x_2]\) are not coordinates in the ordinary affine sense,
because the triple is not uniquely determined by the point:
\[
[x_0:x_1:x_2]
=
[\lambda x_0:\lambda x_1:\lambda x_2]
\qquad
(\lambda\in\Bbbk^\times).
\]
Homogeneous coordinates are therefore redundant labels for projective
points rather than genuine globally defined coordinate functions. By
contrast, once a chart \(U_i\) has been chosen and \(x_i\) normalized
to \(1\), the resulting ratios are honest affine coordinates on that
chart.

In this section, monomials will provide the bridge between these two
descriptions.

The reward will be our first compact polyhedral picture of a genuinely
global object: the global sections of \(\cO_{\PP^2}(d)\) will be
indexed by the lattice points in a triangle of side length \(d\).

\subsection{Why homogeneous polynomials are not ordinary functions}

Let
\[
F(x_0,x_1,x_2)
\]
be a homogeneous polynomial of degree \(d\). If we replace a
representative \((x_0,x_1,x_2)\) of a projective point by the scalar
multiple
\[
(\lambda x_0,\lambda x_1,\lambda x_2),
\]
then
\[
F(\lambda x_0,\lambda x_1,\lambda x_2)
=
\lambda^dF(x_0,x_1,x_2).
\]
Unless \(d=0\), the value of \(F\) therefore depends on the chosen
representative. A homogeneous polynomial of positive degree is not an
ordinary function on \(\PP^2\).

The failure is not arbitrary: changing the representative always
changes the value by the same power \(\lambda^d\). Rather than
discarding homogeneous polynomials, we package this predictable
scaling into a geometric object called the line bundle
\(\cO_{\PP^2}(d)\).

There is an intuitive way to see what this means. On the chart
\(U_0=\{x_0\neq0\}\), divide by \(x_0^d\):
\[
\frac{F(x_0,x_1,x_2)}{x_0^d}
=
F\left(1,\frac{x_1}{x_0},\frac{x_2}{x_0}\right).
\]
This is an honest polynomial function in the affine coordinates
\[
z_1=\frac{x_1}{x_0},
\qquad
z_2=\frac{x_2}{x_0}.
\]
Likewise, on \(U_1\) and \(U_2\), the expressions
\[
\frac{F}{x_1^d}
\qquad\text{and}\qquad
\frac{F}{x_2^d}
\]
are regular functions.

These three local functions do not agree literally on the overlaps.
Instead, they differ by the predictable factors obtained by changing
which homogeneous coordinate has been normalized to \(1\). For
example, on \(U_0\cap U_1\),
\[
\frac{F}{x_0^d}
=
\left(\frac{x_1}{x_0}\right)^d
\frac{F}{x_1^d}.
\]
This controlled compatibility is precisely what it means for \(F\) to
define a global section of \(\cO_{\PP^2}(d)\).

\begin{figure}[ht]
\centering
\begin{tikzpicture}[
  localbox/.style={
    rounded corners,
    minimum width=3.25cm,
    minimum height=1.35cm,
    align=center,
    draw=black!55
  },
  relationbox/.style={
    rounded corners,
    minimum width=8.7cm,
    minimum height=1.05cm,
    align=center,
    draw=ExampleGold!75!black,
    fill=ExampleGold!7
  }
]
  \node[localbox,fill=PastelYellow] (U0) at (-4.0,0.7)
    {\(\displaystyle U_0\)\\[1mm]
     \(\displaystyle F_0=\frac{F}{x_0^d}
       =F(1,z_1,z_2)\)};

  \node[localbox,fill=DictionaryBlue!8] (U1) at (0,0.7)
    {\(\displaystyle U_1\)\\[1mm]
     \(\displaystyle F_1=\frac{F}{x_1^d}\)};

  \node[localbox,fill=SurfaceGreen!9] (U2) at (4.0,0.7)
    {\(\displaystyle U_2\)\\[1mm]
     \(\displaystyle F_2=\frac{F}{x_2^d}\)};

  \node[relationbox] at (0,-1.15)
    {\(\displaystyle
      F_0=\left(\frac{x_1}{x_0}\right)^dF_1
      \quad\text{on }U_0\cap U_1,
      \qquad
      F_1=\left(\frac{x_2}{x_1}\right)^dF_2
      \quad\text{on }U_1\cap U_2.\)};
\end{tikzpicture}
\caption{A homogeneous polynomial of degree \(d\) becomes an ordinary
function after one coordinate is normalized. The lower box records
how the local functions agree on overlaps.}
\label{fig:p2-local-sections}
\end{figure}

\begin{observation}
The line bundle \(\cO_{\PP^2}(d)\) is the device that remembers these
transition factors. No general theory of line bundles is needed at
this stage. For this chapter, one may safely think of a global section
of \(\cO_{\PP^2}(d)\) as a homogeneous polynomial of degree \(d\),
viewed through its compatible local expressions on the three affine
charts.
\end{observation}

For \(d\geq0\), we therefore have
\[
\Hzero{\PP^2}{\cO_{\PP^2}(d)}
\simeq
\Bbbk[x_0,x_1,x_2]_d,
\]
where the right-hand side denotes the vector space of homogeneous
polynomials of degree \(d\).

\subsection{The monomial basis}

A homogeneous monomial of degree \(d\) has the form
\[
x_0^{a_0}x_1^{a_1}x_2^{a_2},
\]
where
\[
a_0,a_1,a_2\in\ZZ_{\geq0},
\qquad
a_0+a_1+a_2=d.
\]
Thus the monomial basis of
\(\Hzero{\PP^2}{\cO_{\PP^2}(d)}\) is indexed by the nonnegative
integer solutions of one linear equation.

There are three exponents but only two independent choices. Set
\[
a=a_1,
\qquad
b=a_2.
\]
Then
\[
a_0=d-a-b,
\]
and the conditions \(a_0,a_1,a_2\geq0\) become
\[
a\geq0,
\qquad
b\geq0,
\qquad
a+b\leq d.
\]
The monomial may therefore be written as
\[
x_0^{d-a-b}x_1^ax_2^b.
\]

On the chart \(U_0\), after dividing by \(x_0^d\), this section becomes
\[
\frac{x_0^{d-a-b}x_1^ax_2^b}{x_0^d}
=
\left(\frac{x_1}{x_0}\right)^a
\left(\frac{x_2}{x_0}\right)^b
=
z_1^az_2^b.
\]
This formula is the key point: a homogeneous monomial on projective
space becomes an ordinary monomial on the affine chart.

\begin{guidingquestion}
Which monomials \(z_1^az_2^b\) can arise from degree-\(d\)
homogeneous monomials?
\end{guidingquestion}

The answer is encoded by the three inequalities
\[
a\geq0,
\qquad
b\geq0,
\qquad
a+b\leq d.
\]
These are exactly the inequalities cutting out a triangle.

\subsection{The triangle of degree-\texorpdfstring{\(d\)}{d} monomials}

Let
\[
\Delta
=
\Conv\{(0,0),(1,0),(0,1)\}
\subseteq M_{\RR}.
\]
Its \(d\)-fold dilation is
\[
d\Delta
=
\Conv\{(0,0),(d,0),(0,d)\}.
\]
The lattice points in \(d\Delta\) are precisely the pairs
\[
(a,b)\in\ZZ^2
\]
satisfying
\[
a\geq0,
\qquad
b\geq0,
\qquad
a+b\leq d.
\]
Hence we obtain the correspondence
\[
x_0^{d-a-b}x_1^ax_2^b
\longleftrightarrow
(a,b)\in d\Delta\cap M.
\]

\begin{figure}[ht]
\centering
\begin{tikzpicture}[scale=0.92]
  \begin{scope}[xshift=-5.0cm]
    \fill[PastelYellow] (0,0) -- (2.2,0) -- (0,2.2) -- cycle;
    \draw[thick] (0,0) -- (2.2,0) -- (0,2.2) -- cycle;
    \foreach \p in {(0,0),(2.2,0),(0,2.2)}
      \fill[DictionaryBlue] \p circle (2pt);
    \node[below] at (1.1,-0.25) {\(d=1\)};
  \end{scope}

  \begin{scope}[xshift=-1.1cm]
    \fill[DictionaryBlue!7] (0,0) -- (2.8,0) -- (0,2.8) -- cycle;
    \draw[thick] (0,0) -- (2.8,0) -- (0,2.8) -- cycle;
    \foreach \i in {0,1,2}
      \foreach \j in {0,...,\numexpr2-\i\relax}
        \fill[ExampleGold] ({1.4*\i},{1.4*\j}) circle (1.8pt);
    \node[below] at (1.4,-0.25) {\(d=2\)};
  \end{scope}

  \begin{scope}[xshift=3.6cm]
    \fill[SurfaceGreen!8] (0,0) -- (3.3,0) -- (0,3.3) -- cycle;
    \draw[thick] (0,0) -- (3.3,0) -- (0,3.3) -- cycle;
    \foreach \i in {0,1,2,3}
      \foreach \j in {0,...,\numexpr3-\i\relax}
        \fill[SurfaceGreen] ({1.1*\i},{1.1*\j}) circle (1.7pt);
    \node[below] at (1.65,-0.25) {\(d=3\)};
  \end{scope}
\end{tikzpicture}
\caption{The lattice points of \(d\Delta\) index the degree-\(d\)
monomials. Increasing \(d\) dilates the triangle and creates more
global sections.}
\label{fig:p2-dilated-triangles}
\end{figure}

For \(d=2\), the six lattice points correspond to the six monomials
\[
x_0^2,\quad
x_0x_1,\quad
x_1^2,\quad
x_0x_2,\quad
x_1x_2,\quad
x_2^2.
\]

\begin{figure}[ht]
\centering
\begin{tikzpicture}[scale=1.25]
  \coordinate (A) at (0,0);
  \coordinate (B) at (4,0);
  \coordinate (C) at (0,4);

  \fill[PastelYellow] (A) -- (B) -- (C) -- cycle;
  \draw[very thick,SurfaceGreen] (A) -- (B);
  \draw[very thick,WarningRed] (B) -- (C);
  \draw[very thick,ExampleGold] (C) -- (A);

  \fill[DictionaryBlue] (0,0) circle (1.9pt);
  \node[below left] at (0,0) {\(x_0^2\)};

  \fill[DictionaryBlue] (2,0) circle (1.9pt);
  \node[below] at (2,0) {\(x_0x_1\)};

  \fill[DictionaryBlue] (4,0) circle (1.9pt);
  \node[below right] at (4,0) {\(x_1^2\)};

  \fill[DictionaryBlue] (0,2) circle (1.9pt);
  \node[left] at (0,2) {\(x_0x_2\)};

  \fill[DictionaryBlue] (2,2) circle (1.9pt);
  \node[above right] at (2,2) {\(x_1x_2\)};

  \fill[DictionaryBlue] (0,4) circle (1.9pt);
  \node[above left] at (0,4) {\(x_2^2\)};

  % The three colored edges are explained immediately below the figure.

\end{tikzpicture}
\caption{The degree-two triangle. The six lattice points label the six
quadratic monomials. The three edge colors match the three coordinate
lines \(D_0,D_1,D_2\).}
\label{fig:p2-degree-two-monomials}
\end{figure}

The colors on the three edges match the coordinate triangle from the
previous sections:
\[
\begin{array}{c|c|c}
\text{edge of }d\Delta
&
\text{missing factor}
&
\text{coordinate line}
\\
\hline
a+b=d & x_0 & D_0=\{x_0=0\}\\
a=0 & x_1 & D_1=\{x_1=0\}\\
b=0 & x_2 & D_2=\{x_2=0\}.
\end{array}
\]
For example, a lattice point on the edge \(a+b=d\) labels a monomial
with exponent \(a_0=0\), so that monomial is not divisible by \(x_0\).
A lattice point in the interior has
\[
a_0>0,\qquad a_1>0,\qquad a_2>0,
\]
so the corresponding monomial is divisible by all three homogeneous
coordinates.

This is our first glimpse of a general principle: the faces of a
polytope record how sections interact with the boundary of a toric
variety.

\subsection{Counting sections by counting lattice points}

The triangle also computes the dimension of the space of global
sections. In the horizontal row with fixed exponent \(b\), the
possible values of \(a\) are
\[
0,1,\ldots,d-b,
\]
so that row contains \(d-b+1\) lattice points. Summing over
\(b=0,\ldots,d\), we find
\[
\#(d\Delta\cap M)
=
1+2+\cdots+(d+1)
=
\frac{(d+1)(d+2)}{2}.
\]
Therefore
\[
\dim_{\Bbbk}
\Hzero{\PP^2}{\cO_{\PP^2}(d)}
=
\frac{(d+1)(d+2)}{2}
=
\binom{d+2}{2}.
\]

This familiar dimension formula is now visible in a picture: it is
simply the number of lattice points in a dilated triangle.

\begin{observation}
Under the torus action
\[
(t_1,t_2)\cdot[x_0:x_1:x_2]
=
[x_0:t_1x_1:t_2x_2],
\]
the monomial section
\[
x_0^{d-a-b}x_1^ax_2^b
\]
is rescaled by
\[
t_1^at_2^b.
\]
In representation-theoretic language, it is a torus
\emph{eigensection} of weight \((a,b)\). No representation theory is
needed to use the correspondence: the lattice point \((a,b)\) simply
records how the monomial changes under coordinatewise rescaling.
\end{observation}

\begin{threedictionaries}
For the line bundle \(\cO_{\PP^2}(d)\), the three viewpoints meet in
one picture.
\begin{center}
\small
\begin{tabularx}{\linewidth}{
  >{\centering\arraybackslash}X|
  >{\centering\arraybackslash}X|
  >{\centering\arraybackslash}X}
\textbf{Algebra} & \textbf{Geometry} & \textbf{Combinatorics}\\
\hline
\(\Bbbk[x_0,x_1,x_2]_d\) &
\(\Hzero{\PP^2}{\cO_{\PP^2}(d)}\) &
\(d\Delta\cap M\)
\\[2mm]
\(x_0^{d-a-b}x_1^ax_2^b\) &
a monomial section &
\((a,b)\)
\\[2mm]
absence of the factor \(x_i\) &
restriction not identically zero on \(D_i\) &
a lattice point on the corresponding edge.
\end{tabularx}
\end{center}
\end{threedictionaries}

\begin{guidingquestion}
The vertices, edges, and interior of the triangle distinguish different
kinds of monomial sections. Do the vertices, edges, and interior also
know something about the torus orbits in \(\PP^2\)?
\end{guidingquestion}


\section{Orbits and coordinate strata}
\label{sec:p2-orbits}
\index{torus orbit}
\index{fixed point!torus-fixed}
\index{orbit closure}

The torus action
\[
(t_1,t_2)\cdot[x_0:x_1:x_2]
=
[x_0:t_1x_1:t_2x_2]
\]
has a simple but powerful feature: it can rescale a nonzero coordinate,
but it can never turn a zero coordinate into a nonzero one, or vice
versa. Thus the pattern of vanishing coordinates is constant along
every orbit.

For \(\PP^2\), there are only seven possible nonempty vanishing
patterns:
\[
\begin{array}{c}
\text{no coordinate vanishes},\\
\text{exactly one coordinate vanishes},\\
\text{exactly two coordinates vanish}.
\end{array}
\]
In fact, each such pattern determines exactly one torus orbit. This
gives a decomposition of \(\PP^2\) into seven pieces.

\begin{guidingquestion}
Why should the vanishing pattern determine the entire orbit, rather
than merely a union of several orbits?
\end{guidingquestion}

The reason is that the torus has enough freedom to change all the
nonzero coordinates relative to one another. Let us work through the
three cases.

\subsection{The open orbit}

If no coordinate vanishes, the point belongs to
\[
T
=
\{[x_0:x_1:x_2]:x_0x_1x_2\neq0\}.
\]
We already know that \(T\) is a single orbit. Indeed, after
normalizing \(x_0=1\), any point has the form
\[
[1:z_1:z_2],
\qquad
z_1,z_2\in\Bbbk^\times,
\]
and
\[
(z_1^{-1},z_2^{-1})\cdot[1:z_1:z_2]
=
[1:1:1].
\]
Thus every point of \(T\) can be moved to the identity
\([1:1:1]\).

This orbit has dimension \(2\), the same as \(\PP^2\), and it is open
and dense.

\subsection{The three one-dimensional orbits}

Suppose exactly one homogeneous coordinate vanishes. For example,
consider the coordinate line
\[
D_0=\{x_0=0\}.
\]
Removing its two coordinate points leaves
\[
O_0
=
D_0\setminus\{p_1,p_2\}
=
\{[0:x_1:x_2]:x_1x_2\neq0\}.
\]
Every point of \(O_0\) may be written uniquely as
\[
[0:1:z],
\qquad
z\in\Bbbk^\times.
\]
The torus action becomes
\[
(t_1,t_2)\cdot[0:1:z]
=
[0:t_1:t_2z]
=
\left[0:1:\frac{t_2}{t_1}z\right].
\]
By choosing the ratio \(t_2/t_1\), we may move any point of \(O_0\)
to any other. Hence \(O_0\) is a single orbit.

Exactly the same argument applies to the other two coordinate lines:
\[
\begin{aligned}
O_1
&=
D_1\setminus\{p_0,p_2\}
=
\{[x_0:0:x_2]:x_0x_2\neq0\},
\\
O_2
&=
D_2\setminus\{p_0,p_1\}
=
\{[x_0:x_1:0]:x_0x_1\neq0\}.
\end{aligned}
\]
Each \(O_i\) is isomorphic to \(\Gm\), and each is a
one-dimensional torus orbit.

\begin{observation}
The torus \(T\simeq\Gm^2\) acts transitively on \(O_i\simeq\Gm\), but
not freely: one direction of the two-dimensional torus acts
trivially. For instance, on \(O_0\), only the ratio \(t_2/t_1\)
matters. The diagonal subgroup
\[
\{(t,t):t\in\Gm\}
\]
fixes every point of \(O_0\). This explains geometrically why the
orbit has dimension \(1\) rather than \(2\).
\end{observation}

\subsection{The three fixed points}

If exactly two coordinates vanish, the point must be one of
\[
p_0=[1:0:0],
\qquad
p_1=[0:1:0],
\qquad
p_2=[0:0:1].
\]
Each is fixed by the entire torus. For example,
\[
(t_1,t_2)\cdot[1:0:0]
=
[1:0:0].
\]
Thus each coordinate point is a zero-dimensional orbit.

We have proved the orbit decomposition
\[
\PP^2
=
T
\sqcup O_0\sqcup O_1\sqcup O_2
\sqcup\{p_0\}\sqcup\{p_1\}\sqcup\{p_2\}.
\]

\begin{figure}[ht]
\centering
\begin{tikzpicture}[scale=1.05]
  \coordinate (p0) at (0,3.4);
  \coordinate (p1) at (-3.2,0);
  \coordinate (p2) at (3.2,0);

  % Open orbit.
  \fill[PastelYellow]
    (0,2.75) -- (-2.35,0.42) -- (2.35,0.42) -- cycle;
  \node[align=center] at (0,1.45)
    {\(\displaystyle T\)\\[-1mm]
     \(2\)-dimensional orbit};

  % One-dimensional orbits.
  \draw[very thick,SurfaceGreen] (p0) -- (p1);
  \draw[very thick,ExampleGold] (p0) -- (p2);
  \draw[very thick,WarningRed] (p1) -- (p2);

  \node[rotate=47,fill=white,rounded corners,inner sep=2pt,
        text=SurfaceGreen!65!black]
    at (-1.55,1.75) {\(O_2\)};
  \node[rotate=-47,fill=white,rounded corners,inner sep=2pt,
        text=ExampleGold!75!black]
    at (1.55,1.75) {\(O_1\)};
  \node[fill=white,rounded corners,inner sep=2pt,text=WarningRed]
    at (0,0) {\(O_0\)};

  % Fixed points.
  \fill[DictionaryBlue] (p0) circle (2.4pt);
  \fill[DictionaryBlue] (p1) circle (2.4pt);
  \fill[DictionaryBlue] (p2) circle (2.4pt);

  \node[above] at (p0) {\(p_0\)};
  \node[below left] at (p1) {\(p_1\)};
  \node[below right] at (p2) {\(p_2\)};
\end{tikzpicture}
\caption{The seven torus orbits of \(\PP^2\): one open
two-dimensional orbit, three one-dimensional orbits, and three fixed
points. As before, the triangular drawing is a combinatorial model of
the coordinate-line incidence pattern.}
\label{fig:p2-seven-orbits}
\end{figure}

\subsection{Orbit closures}

The orbits do not sit independently inside \(\PP^2\). Their closures
remember how one stratum can degenerate into another.

The closure of the open orbit is all of \(\PP^2\):
\[
\overline{T}=\PP^2.
\]
The closure of a one-dimensional orbit is the corresponding coordinate
line:
\[
\overline{O_i}=D_i.
\]
For example,
\[
\overline{O_0}
=
D_0
=
O_0\sqcup\{p_1,p_2\}.
\]
Finally, each fixed point is already closed:
\[
\overline{\{p_i\}}=\{p_i\}.
\]

These relations may be displayed in a small incidence diagram.

\begin{figure}[ht]
\centering
\begin{tikzpicture}[
  orbit/.style={
    rounded corners,
    draw=black!55,
    minimum width=1.35cm,
    minimum height=0.75cm,
    align=center
  }
]
  \node[orbit,fill=PastelYellow] (T) at (0,3.2) {\(T\)};

  \node[orbit,fill=WarningRed!7] (O0) at (-3,1.55) {\(O_0\)};
  \node[orbit,fill=ExampleGold!8] (O1) at (0,1.55) {\(O_1\)};
  \node[orbit,fill=SurfaceGreen!8] (O2) at (3,1.55) {\(O_2\)};

  \node[orbit,fill=DictionaryBlue!8] (p0) at (-3,-0.15) {\(p_0\)};
  \node[orbit,fill=DictionaryBlue!8] (p1) at (0,-0.15) {\(p_1\)};
  \node[orbit,fill=DictionaryBlue!8] (p2) at (3,-0.15) {\(p_2\)};

  \draw[thick] (T) -- (O0);
  \draw[thick] (T) -- (O1);
  \draw[thick] (T) -- (O2);

  % O0 closure contains p1,p2
  \draw[thick,WarningRed] (O0) -- (p1);
  \draw[thick,WarningRed] (O0) -- (p2);

  % O1 closure contains p0,p2
  \draw[thick,ExampleGold] (O1) -- (p0);
  \draw[thick,ExampleGold] (O1) -- (p2);

  % O2 closure contains p0,p1
  \draw[thick,SurfaceGreen] (O2) -- (p0);
  \draw[thick,SurfaceGreen] (O2) -- (p1);
\end{tikzpicture}
\caption{The orbit-closure relations. A downward path means that the
lower orbit lies in the closure of the upper one.}
\label{fig:p2-orbit-closure-poset}
\end{figure}

There is also a direct way to see these degenerations, but in
projective space one must remember that a limit may become visible
only after changing the homogeneous representative.

For \((u,v)\in\ZZ^2\), consider the one-parameter motion
\[
\lambda^{(u,v)}(t)=(t^u,t^v)\in T.
\]
Acting on the identity point \(e=[1:1:1]\) gives
\[
\gamma_{u,v}(t)
=
\lambda^{(u,v)}(t)\cdot e
=
[1:t^u:t^v].
\]
The exponents \(u\) and \(v\) are allowed to be negative. Although an
individual coordinate may then have a pole at \(t=0\), the
\emph{projective point} can still have a limit, because we may rescale
all three homogeneous coordinates by the same power of \(t\).

Set
\[
m=\min\{0,u,v\}.
\]
For \(t\neq0\), we may rewrite the same projective point as
\[
\gamma_{u,v}(t)
=
[t^{-m}:t^{u-m}:t^{v-m}].
\]
All three exponents are now nonnegative, and at least one is zero.
Consequently, the limit is obtained by keeping precisely those
coordinates whose original exponent attains the minimum of
\(\{0,u,v\}\).

This rule reaches strata that are invisible in the affine chart
\(U_0\). For example,
\[
\begin{aligned}
\lambda^{(1,1)}(t)\cdot e
&=[1:t:t]
  \longrightarrow [1:0:0]=p_0,
\\
\lambda^{(-1,0)}(t)\cdot e
&=[1:t^{-1}:1]
 =[t:1:t]
  \longrightarrow [0:1:0]=p_1,
\\
\lambda^{(0,-1)}(t)\cdot e
&=[1:1:t^{-1}]
 =[t:t:1]
  \longrightarrow [0:0:1]=p_2.
\end{aligned}
\]
The second and third limits land naturally in the charts \(U_1\) and
\(U_2\), respectively.

The same idea reaches distinguished points on the
one-dimensional orbits:
\[
\begin{aligned}
\lambda^{(-1,-1)}(t)\cdot e
&=[1:t^{-1}:t^{-1}]
 =[t:1:1]
  \longrightarrow [0:1:1]\in O_0,
\\
\lambda^{(1,0)}(t)\cdot e
&=[1:t:1]
  \longrightarrow [1:0:1]\in O_1,
\\
\lambda^{(0,1)}(t)\cdot e
&=[1:1:t]
  \longrightarrow [1:1:0]\in O_2.
\end{aligned}
\]

More generally, begin at an arbitrary point
\[
q=[1:c_1:c_2]\in T.
\]
If \(u=v<0\), then
\[
\lambda^{(u,u)}(t)\cdot q
=
[1:c_1t^u:c_2t^u]
=
[t^{-u}:c_1:c_2]
\longrightarrow
[0:c_1:c_2]\in O_0.
\]
As \(c_1\) and \(c_2\) vary, every point of \(O_0\) occurs in this
way. Analogous choices produce every point of \(O_1\) and \(O_2\).

Thus the limit is governed not by whether \(u\) and \(v\) are
nonnegative in one chosen chart, but by which of the three numbers
\[
0,\qquad u,\qquad v
\]
are minimal. A unique minimum produces a fixed point; a minimum
attained twice produces a one-dimensional orbit; and equality of all
three exponents leaves us in the open orbit. This is the projective
counterpart of the limit behavior we saw in \(\AAA^2\).

\subsection{Two triangles, two different meanings}

At this point, two triangular objects have appeared:

\begin{enumerate}
\item the \emph{coordinate triangle}
\[
D_0\cup D_1\cup D_2\subseteq\PP^2,
\]
which is the union of three projective lines;

\item the \emph{lattice triangle}
\[
\Delta
=
\Conv\{(0,0),(1,0),(0,1)\}
\subseteq M_{\RR},
\]
which is an actual Euclidean triangle in a real vector space.
\end{enumerate}

These should not be confused. The first is a reducible algebraic curve
inside \(\PP^2\); the second is a convex polygon in \(M_{\RR}\).
Nevertheless, they encode the same incidence pattern, and this is one
of the central miracles of toric geometry.

The face structure of \(\Delta\) matches the orbit structure of
\(\PP^2\):
\[
\begin{array}{c|c|c}
\text{face of }\Delta
&
\text{torus orbit in }\PP^2
&
\text{dimension}
\\
\hline
\Delta
&
T
&
2
\\
\text{an edge}
&
O_i
&
1
\\
\text{a vertex}
&
p_i
&
0.
\end{array}
\]
Here the whole triangle \(\Delta\), regarded as its unique
two-dimensional face, corresponds to the open orbit. The edges
correspond to the punctured coordinate lines, and the vertices
correspond to the fixed points.

To make the bijection completely explicit, regard \(\Delta\) as a
cell complex whose cells are the relative interiors of its faces.

\begin{center}
\small
\begin{tabularx}{\linewidth}{@{}
  >{\raggedright\arraybackslash}X
  >{\raggedright\arraybackslash}X
  >{\raggedright\arraybackslash}X
  >{\raggedright\arraybackslash}X@{}}
\toprule
\textbf{Cell of \(\Delta\)}
&
\textbf{Closed face}
&
\textbf{Orbit}
&
\textbf{Representative}
\\
\midrule
\(\relint(\Delta)\)
&
\(\Delta\)
&
\(T\)
&
\([1:1:1]\)
\\[1mm]
\(\{a,b>0,\ a+b=1\}\)
&
\(F_0=\{a+b=1\}\)
&
\(O_0\)
&
\([0:1:1]\)
\\[1mm]
\(\{a=0,\ 0<b<1\}\)
&
\(F_1=\{a=0\}\)
&
\(O_1\)
&
\([1:0:1]\)
\\[1mm]
\(\{b=0,\ 0<a<1\}\)
&
\(F_2=\{b=0\}\)
&
\(O_2\)
&
\([1:1:0]\)
\\[1mm]
\(\{(0,0)\}\)
&
\(\{(0,0)\}\)
&
\(\{p_0\}\)
&
\([1:0:0]\)
\\[1mm]
\(\{(1,0)\}\)
&
\(\{(1,0)\}\)
&
\(\{p_1\}\)
&
\([0:1:0]\)
\\[1mm]
\(\{(0,1)\}\)
&
\(\{(0,1)\}\)
&
\(\{p_2\}\)
&
\([0:0:1]\)
\\
\bottomrule
\end{tabularx}
\end{center}

Passing from a cell to its closed face corresponds geometrically to
passing from an orbit to its orbit closure:
\[
F_0\longleftrightarrow D_0,\qquad
F_1\longleftrightarrow D_1,\qquad
F_2\longleftrightarrow D_2.
\]

\begin{figure}[ht]
\centering
\begin{tikzpicture}[scale=0.92]
  % Coordinate triangle/orbits.
  \begin{scope}[xshift=-4.6cm]
    \coordinate (P0) at (0,3.0);
    \coordinate (P1) at (-2.7,0);
    \coordinate (P2) at (2.7,0);

    \fill[PastelYellow] (0,2.35) -- (-1.95,0.38) -- (1.95,0.38) -- cycle;
    \draw[very thick,SurfaceGreen] (P0) -- (P1);
    \draw[very thick,ExampleGold] (P0) -- (P2);
    \draw[very thick,WarningRed] (P1) -- (P2);

    \fill[DictionaryBlue] (P0) circle (2.2pt);
    \fill[DictionaryBlue] (P1) circle (2.2pt);
    \fill[DictionaryBlue] (P2) circle (2.2pt);

    \node at (0,1.25) {\(T\)};
    \node[below] at (0,-0.45) {orbits in \(\PP^2\)};
  \end{scope}

  % Arrow.
  \draw[very thick,-{Latex[length=2.6mm]},DictionaryBlue]
    (-1.25,1.45) -- (1.25,1.45)
    node[midway,above] {same incidence pattern};

  % Lattice triangle.
  \begin{scope}[xshift=4.6cm]
    \fill[PastelYellow] (0,0) -- (3.0,0) -- (0,3.0) -- cycle;
    \draw[very thick,SurfaceGreen] (0,0) -- (3.0,0);
    \draw[very thick,WarningRed] (3.0,0) -- (0,3.0);
    \draw[very thick,ExampleGold] (0,3.0) -- (0,0);

    \fill[DictionaryBlue] (0,0) circle (2.2pt);
    \fill[DictionaryBlue] (3.0,0) circle (2.2pt);
    \fill[DictionaryBlue] (0,3.0) circle (2.2pt);

    \node at (0.9,0.9) {\(\Delta\)};
    \node[below] at (1.5,-0.45) {faces of \(\Delta\)};
  \end{scope}
\end{tikzpicture}
\caption{The coordinate triangle in \(\PP^2\) and the lattice triangle
\(\Delta\) are different mathematical objects, but their strata have
the same incidence pattern. Matching colors indicate the
corresponding one-dimensional strata.}
\label{fig:p2-orbits-faces}
\end{figure}

\begin{warningbox}
A lattice point in the interior of \(d\Delta\) does not by itself
represent a torus orbit. Individual lattice points index monomial
sections. Torus orbits correspond to faces of the polytope: the whole
polytope, its edges, and its vertices.
\end{warningbox}

\begin{threedictionaries}
The orbit decomposition of \(\PP^2\) can be read in three parallel
languages.

\medskip

\noindent
\begin{minipage}[t]{0.30\linewidth}
\centering
\textbf{Algebra}\par
\smallskip
Record which of
\[
x_0,\ x_1,\ x_2
\]
vanish. The vanishing set is constant along an orbit.
\end{minipage}
\hfill
\begin{minipage}[t]{0.30\linewidth}
\centering
\textbf{Geometry}\par
\smallskip
The vanishing pattern determines one torus orbit:
\[
T,\quad O_i,\quad\text{or}\quad p_i.
\]
Allowing additional coordinates to vanish describes orbit closure.
\end{minipage}
\hfill
\begin{minipage}[t]{0.30\linewidth}
\centering
\textbf{Combinatorics}\par
\smallskip
The orbit corresponds to the relative interior of a face of
\(\Delta\). Orbit closure corresponds to replacing that cell by its
closed face.
\end{minipage}

\medskip

\[
\begin{array}{c@{\quad\longleftrightarrow\quad}c
              @{\quad\longleftrightarrow\quad}c}
\text{vanishing pattern}
&
\text{torus orbit}
&
\text{cell of }\Delta
\end{array}
\]
\end{threedictionaries}

\begin{guidingquestion}
The triangle \(\Delta\) explains the orbit structure beautifully, but
it lives in the space of monomial exponents. Is there another
combinatorial object that records how the three affine charts of
\(\PP^2\) fit together?
\end{guidingquestion}


\section{The triangle and the fan}
\label{sec:p2-triangle-fan}
\index{fan!of the projective plane}
\index{normal fan}
\index{one-parameter subgroup!limits}

The lattice triangle
\[
\Delta=\Conv\{(0,0),(1,0),(0,1)\}\subseteq M_{\RR}
\]
has already encoded two global features of \(\PP^2\): its monomial
sections and its torus orbits. We now extract a second combinatorial
object from the same triangle.

There is one important change of viewpoint. The triangle lives in
\(M_{\RR}\), the space of monomial exponents, whereas the new object
lives in the paired vector space \(N_{\RR}\), where we record
one-parameter motions. Its job is different as well: rather than
recording sections, it remembers the three affine charts of
\(\PP^2\) and how they fit together.

\subsection{Three sides, three inward directions}

Write a point of \(M_{\RR}\) as \(m=(a,b)\). The triangle \(\Delta\)
is cut out by the inequalities
\[
a\geq0,\qquad b\geq0,\qquad a+b\leq1.
\]
Its three edges are
\[
F_0=\{a+b=1\},\qquad
F_1=\{a=0\},\qquad
F_2=\{b=0\}.
\]

Each edge has a distinguished primitive direction pointing into the
triangle:
\[
v_0=(-1,-1),\qquad
v_1=(1,0),\qquad
v_2=(0,1).
\]
Indeed, moving in the direction \(v_1\) increases \(a\), moving in the
direction \(v_2\) increases \(b\), and moving in the direction \(v_0\)
decreases \(a+b\). Thus \(v_i\) points inward from the edge \(F_i\).

\begin{figure}[ht]
\centering
\begin{tikzpicture}[scale=1.45]
  \coordinate (A) at (0,0);
  \coordinate (B) at (3.2,0);
  \coordinate (C) at (0,3.2);

  \fill[PastelYellow] (A) -- (B) -- (C) -- cycle;
  \draw[very thick,SurfaceGreen] (A) -- (B);
  \draw[very thick,WarningRed] (B) -- (C);
  \draw[very thick,ExampleGold] (C) -- (A);

  \fill[DictionaryBlue] (A) circle (2pt);
  \fill[DictionaryBlue] (B) circle (2pt);
  \fill[DictionaryBlue] (C) circle (2pt);

  \node[below left] at (A) {\((0,0)\)};
  \node[below right] at (B) {\((1,0)\)};
  \node[above left] at (C) {\((0,1)\)};

  \draw[very thick,-{Latex[length=2.5mm]},WarningRed]
    (2.05,1.15) -- (1.50,0.60);
  \node[right,WarningRed] at (2.12,1.28) {\(v_0=(-1,-1)\)};

  \draw[very thick,-{Latex[length=2.5mm]},ExampleGold]
    (0,1.55) -- (0.70,1.55);
  \node[left,ExampleGold!75!black] at (-0.05,1.55) {\(v_1=(1,0)\)};

  \draw[very thick,-{Latex[length=2.5mm]},SurfaceGreen]
    (1.55,0) -- (1.55,0.70);
  \node[below,SurfaceGreen!70!black] at (1.55,-0.08)
    {\(v_2=(0,1)\)};
\end{tikzpicture}
\caption{The three primitive inward normals of the standard triangle.
The colors match the coordinate lines \(D_i\) and the corresponding
edges \(F_i\).}
\label{fig:p2-inward-normals}
\end{figure}

There is a mild abuse hidden in this picture. The triangle lies in
\(M_{\RR}\), while the normal vectors lie in \(N_{\RR}\). After
choosing dual bases, both spaces look like \(\RR^2\), so it is useful
to draw the vectors next to the triangle. Conceptually, however, they
belong to paired spaces.

\subsection{A direction asks a question of the triangle}

Let \(n=(u,v)\in N_{\RR}\). Pairing with \(n\) produces a linear
function on the triangle:
\[
m=(a,b)\longmapsto\pair{m}{n}=au+bv.
\]
We ask:

\begin{center}
\emph{On which part of \(\Delta\) is this linear function smallest?}
\end{center}

The answer depends only on the direction of \(n\).

For example, if \(n=v_1=(1,0)\), then
\[
\pair{(a,b)}{v_1}=a.
\]
Its minimum is attained on the entire edge \(F_1=\{a=0\}\). Similarly,
\(v_2\) singles out \(F_2\), while \(v_0\) singles out \(F_0\).

A more general direction can single out a vertex. If \(u>0\) and
\(v>0\), then \(au+bv\) is smallest uniquely at \((0,0)\). The set of
all such directions, together with its boundary, is the cone
\[
\sigma_0=\Cone{v_1,v_2}.
\]
Likewise,
\[
\sigma_1=\Cone{v_2,v_0}
\]
is the set of directions whose minimum is attained at \((1,0)\), and
\[
\sigma_2=\Cone{v_0,v_1}
\]
is the set of directions whose minimum is attained at \((0,1)\).

\begin{figure}[ht]
\centering
\begin{tikzpicture}[scale=1.2]
  % Colored maximal cones.
  \fill[PastelYellow] (0,0) -- (3.0,0) -- (0,3.0) -- cycle;
  \fill[DictionaryBlue!9] (0,0) -- (0,3.0) -- (-2.7,-2.7) -- cycle;
  \fill[SurfaceGreen!10] (0,0) -- (-2.7,-2.7) -- (3.0,0) -- cycle;

  % Rays.
  \draw[very thick,-{Latex[length=2.8mm]},WarningRed]
    (0,0) -- (-2.9,-2.9);
  \draw[very thick,-{Latex[length=2.8mm]},ExampleGold]
    (0,0) -- (3.25,0);
  \draw[very thick,-{Latex[length=2.8mm]},SurfaceGreen]
    (0,0) -- (0,3.25);

  % Labels for rays.
  \node[below left,WarningRed] at (-2.65,-2.65) {\(\rho_0\)};
  \node[below,ExampleGold!75!black] at (2.65,0) {\(\rho_1\)};
  \node[left,SurfaceGreen!70!black] at (0,2.65) {\(\rho_2\)};

  % Labels for cones.
  \node at (1.15,1.05) {\(\sigma_0\)};
  \node at (-1.05,0.65) {\(\sigma_1\)};
  \node at (0.55,-1.20) {\(\sigma_2\)};

  \fill (0,0) circle (1.7pt);
  \node[above right] at (0.08,0.08) {\(\{0\}\)};
\end{tikzpicture}
\caption{The normal fan of \(\Delta\). The rays are
\(\rho_i=\RR_{\geq0}v_i\), and the three two-dimensional cones are
\(\sigma_0,\sigma_1,\sigma_2\). Together they cover \(N_{\RR}\).}
\label{fig:p2-normal-fan}
\end{figure}

The collection
\[
\Sigma_{\Delta}
=
\bigl\{
\{0\},
\rho_0,\rho_1,\rho_2,
\sigma_0,\sigma_1,\sigma_2
\bigr\}
\]
is called the \emph{normal fan} of \(\Delta\). Here
\[
\rho_i=\RR_{\geq0}v_i.
\]
The word \emph{fan} refers to a collection of cones meeting cleanly
along common faces. In this example, the three two-dimensional cones
fit around the origin like the sectors of a fan and cover the whole
plane \(N_{\RR}\).

There is an intuitive reason they cover everything. Every linear
function on the compact triangle \(\Delta\) attains a minimum
somewhere: at a vertex, along an edge, or on all of \(\Delta\) if the
function is constant. The normal fan records exactly where that
minimum occurs.

\subsection{Faces of the triangle and cones of the fan}

The passage from a face of \(\Delta\) to its normal cone reverses
dimensions:
\[
\begin{array}{c|c}
\text{face of }\Delta & \text{normal cone in }\Sigma_\Delta\\
\hline
\Delta & \{0\}\\
F_i & \rho_i\\
\text{a vertex} & \text{a two-dimensional cone}.
\end{array}
\]
More explicitly,
\[
\begin{array}{c|c}
\text{face of }\Delta & \text{normal cone}\\
\hline
\Delta & \{0\}\\
F_0=\{a+b=1\} & \rho_0=\RR_{\geq0}v_0\\
F_1=\{a=0\} & \rho_1=\RR_{\geq0}v_1\\
F_2=\{b=0\} & \rho_2=\RR_{\geq0}v_2\\
(0,0) & \sigma_0=\Cone{v_1,v_2}\\
(1,0) & \sigma_1=\Cone{v_2,v_0}\\
(0,1) & \sigma_2=\Cone{v_0,v_1}.
\end{array}
\]

If one face is contained in another, the corresponding normal cones
are contained in the opposite order. For example,
\[
(0,0)\in F_1
\qquad\Longleftrightarrow\qquad
\rho_1\subseteq\sigma_0.
\]
This reversal is not an accident. Requiring a linear function to be
minimal along an entire edge is more restrictive than requiring it to
be minimal at one of the edge's vertices. The larger face therefore
has the smaller normal cone.

\subsection{The fan and one-parameter limits}

The fan packages the limit calculations of the previous section.

A direction in the relative interior of a ray \(\rho_i\) leads to the
one-dimensional orbit \(O_i\). For the primitive ray generators,
\[
\begin{aligned}
v_0=(-1,-1)
&:\quad
[1:t^{-1}:t^{-1}]
=
[t:1:1]
\longrightarrow[0:1:1]\in O_0,
\\
v_1=(1,0)
&:\quad
[1:t:1]
\longrightarrow[1:0:1]\in O_1,
\\
v_2=(0,1)
&:\quad
[1:1:t]
\longrightarrow[1:1:0]\in O_2.
\end{aligned}
\]

A direction in the relative interior of a two-dimensional cone leads
to the corresponding fixed point:
\[
\begin{array}{c|c}
\text{cone containing }(u,v) & 
\displaystyle\lim_{t\to0}[1:t^u:t^v]\\
\hline
\relint(\sigma_0) & p_0=[1:0:0]\\
\relint(\sigma_1) & p_1=[0:1:0]\\
\relint(\sigma_2) & p_2=[0:0:1].
\end{array}
\]
Finally, the zero vector gives the constant point \([1:1:1]\) in the
open orbit.

Thus cones of dimension \(k\) correspond to orbits of dimension
\(2-k\):
\[
\dim O(\sigma)=2-\dim\sigma.
\]
At this stage, the formula is a summary of the seven calculations
above. Later it will become part of the general orbit--cone
correspondence.

\begin{figure}[ht]
\centering
\begin{tikzpicture}[scale=0.9]
  % Fan on left.
  \begin{scope}[xshift=-4.5cm]
    \fill[PastelYellow] (0,0) -- (2.6,0) -- (0,2.6) -- cycle;
    \fill[DictionaryBlue!9] (0,0) -- (0,2.6) -- (-2.25,-2.25) -- cycle;
    \fill[SurfaceGreen!10] (0,0) -- (-2.25,-2.25) -- (2.6,0) -- cycle;

    \draw[very thick,WarningRed] (0,0) -- (-2.45,-2.45);
    \draw[very thick,ExampleGold] (0,0) -- (2.8,0);
    \draw[very thick,SurfaceGreen] (0,0) -- (0,2.8);
    \fill (0,0) circle (1.5pt);

    \node at (0.95,0.9) {\(\sigma_0\)};
    \node at (-0.9,0.55) {\(\sigma_1\)};
    \node at (0.45,-1.0) {\(\sigma_2\)};
    \node[below] at (0,-2.9) {cones in \(N_{\RR}\)};
  \end{scope}

  \draw[very thick,-{Latex[length=2.6mm]},DictionaryBlue]
    (-1.1,0.25) -- (1.1,0.25)
    node[midway,above,align=center]
    {one-parameter\\limits};

  % Orbit triangle on right.
  \begin{scope}[xshift=4.5cm]
    \coordinate (P0) at (0,2.85);
    \coordinate (P1) at (-2.55,0);
    \coordinate (P2) at (2.55,0);

    \fill[PastelYellow] (0,2.20) -- (-1.82,0.35) -- (1.82,0.35) -- cycle;
    \draw[very thick,SurfaceGreen] (P0) -- (P1);
    \draw[very thick,ExampleGold] (P0) -- (P2);
    \draw[very thick,WarningRed] (P1) -- (P2);

    \fill[DictionaryBlue] (P0) circle (2.1pt);
    \fill[DictionaryBlue] (P1) circle (2.1pt);
    \fill[DictionaryBlue] (P2) circle (2.1pt);

    \node[above] at (P0) {\(p_0\)};
    \node[below left] at (P1) {\(p_1\)};
    \node[below right] at (P2) {\(p_2\)};
    \node at (0,1.15) {\(T\)};
    \node[below] at (0,-0.55) {orbits in \(\PP^2\)};
  \end{scope}
\end{tikzpicture}
\caption{The cone containing a one-parameter direction determines the
orbit in which the limit lies. Two-dimensional cones lead to fixed
points, rays lead to one-dimensional orbits, and the zero cone
corresponds to the open orbit.}
\label{fig:p2-fan-orbits}
\end{figure}

\subsection{The three maximal cones recover the three affine charts}

The fan remembers more than the orbit decomposition. Each
two-dimensional cone recovers one of the standard affine charts.

Start with
\[
\sigma_0=\Cone{v_1,v_2},
\]
the first quadrant in \(N_{\RR}\). A monomial \(z_1^az_2^b\) extends
regularly along every one-parameter motion in \(\sigma_0\) precisely
when
\[
a\geq0,\qquad b\geq0.
\]
These monomials form the polynomial ring
\[
\Bbbk[z_1,z_2],
\]
whose spectrum is the chart
\[
U_0\simeq\AAA^2.
\]

For the second cone,
\[
\sigma_1=\Cone{v_2,v_0},
\]
the extendable Laurent monomials are generated by
\[
z_1^{-1}
\qquad\text{and}\qquad
z_1^{-1}z_2.
\]
But these are exactly the affine coordinates on \(U_1\):
\[
w_0=\frac{x_0}{x_1}=z_1^{-1},
\qquad
w_2=\frac{x_2}{x_1}=z_1^{-1}z_2.
\]
Hence
\[
\Bbbk[\sigma_1^\vee\cap M]
=
\Bbbk[w_0,w_2],
\qquad
U_1\simeq\AAA^2.
\]

Similarly, for
\[
\sigma_2=\Cone{v_0,v_1},
\]
the extendable Laurent monomials are generated by
\[
z_2^{-1}
\qquad\text{and}\qquad
z_1z_2^{-1},
\]
which are the coordinates
\[
r_0=\frac{x_0}{x_2},
\qquad
r_1=\frac{x_1}{x_2}
\]
on \(U_2\). Thus
\[
\Bbbk[\sigma_2^\vee\cap M]
=
\Bbbk[r_0,r_1].
\]

\begin{center}
\small
\begin{tabularx}{\linewidth}{@{}
  >{\raggedright\arraybackslash}X
  >{\raggedright\arraybackslash}X
  >{\raggedright\arraybackslash}X
  >{\raggedright\arraybackslash}X@{}}
\toprule
\textbf{Vertex of \(\Delta\)}
&
\textbf{Maximal cone}
&
\textbf{Affine coordinates}
&
\textbf{Chart and fixed point}
\\
\midrule
\((0,0)\)
&
\(\sigma_0=\Cone{v_1,v_2}\)
&
\(z_1=x_1/x_0,\ z_2=x_2/x_0\)
&
\(U_0,\ p_0\)
\\[1mm]
\((1,0)\)
&
\(\sigma_1=\Cone{v_2,v_0}\)
&
\(w_0=x_0/x_1,\ w_2=x_2/x_1\)
&
\(U_1,\ p_1\)
\\[1mm]
\((0,1)\)
&
\(\sigma_2=\Cone{v_0,v_1}\)
&
\(r_0=x_0/x_2,\ r_1=x_1/x_2\)
&
\(U_2,\ p_2\)
\\
\bottomrule
\end{tabularx}
\end{center}

This table is the first complete appearance of our three dictionaries:
a vertex of the polytope determines a cone; the cone determines a
semigroup of monomials; and those monomials form the coordinate ring
of an affine chart.

\subsection{What the fan remembers}

For \(\PP^2\), the normal fan packages the following correspondences:
\[
\begin{array}{c@{\quad\longleftrightarrow\quad}c}
\{0\} & \text{the dense torus }T,\\
\rho_i & \text{the coordinate line }D_i,\\
\sigma_i & \text{the affine chart }U_i
            \text{ and its fixed point }p_i.
\end{array}
\]
Containment of cones reverses containment of orbit closures. For
example,
\[
\rho_1,\rho_2\subseteq\sigma_0
\]
records that
\[
p_0=D_1\cap D_2.
\]

\begin{threedictionaries}
The normal fan translates between local algebra, projective geometry,
and convex combinatorics.

\medskip

\noindent
\begin{minipage}[t]{0.30\linewidth}
\centering
\textbf{Algebra}\par
\smallskip
A cone \(\sigma_i\) selects the Laurent monomials that extend
regularly. They form the coordinate ring
\[
\Bbbk[\sigma_i^\vee\cap M].
\]
\end{minipage}
\hfill
\begin{minipage}[t]{0.30\linewidth}
\centering
\textbf{Geometry}\par
\smallskip
That ring describes the affine chart \(U_i\). Faces of the cone
describe the torus orbits meeting the chart.
\end{minipage}
\hfill
\begin{minipage}[t]{0.30\linewidth}
\centering
\textbf{Combinatorics}\par
\smallskip
Vertices, edges, and the whole polytope have normal cones of
dimensions \(2,1,0\), respectively.
\end{minipage}
\end{threedictionaries}

\begin{surfacelaboratory}
Starting from
\[
\Delta=\Conv\{(0,0),(1,0),(0,1)\},
\]
perform the following reconstruction of \(\PP^2\).

\begin{enumerate}
\item Find the primitive inward normal to each edge of \(\Delta\).
\item Form the three two-dimensional cones generated by adjacent
normal vectors.
\item For each maximal cone, determine the Laurent monomials that pair
nonnegatively with it.
\item Recognize the resulting three polynomial rings as the coordinate
rings of \(U_0,U_1,U_2\).
\end{enumerate}
\end{surfacelaboratory}

\begin{guidingquestion}
If changing the fan changes the way affine charts fit together, what
geometric operation should correspond to inserting a new ray into one
of its two-dimensional cones?
\end{guidingquestion}


\section{Blowing up a torus-fixed point}
\label{sec:p2-first-blowup}
\index{blowup!of the projective plane}
\index{star subdivision}
\index{exceptional divisor}
\index{self-intersection}

The fan of \(\PP^2\) contains three two-dimensional cones, one for
each torus-fixed point. We now modify one of these cones and watch the
projective plane change.

Our chosen point is
\[
p_0=[1:0:0].
\]
It lies in the affine chart
\[
U_0=\{x_0\neq0\}\simeq\AAA^2
\]
and becomes the origin in the affine coordinates
\[
z_1=\frac{x_1}{x_0},
\qquad
z_2=\frac{x_2}{x_0}.
\]
Thus blowing up \(p_0\) in \(\PP^2\) is, near \(p_0\), the familiar
operation of blowing up the origin in the affine plane.

\subsection{Replacing a point by its directions}

The origin of \(\AAA^2\) has many possible directions of approach.
A line through the origin is determined by a projective point
\[
[s:t]\in\PP^1.
\]
The blowup remembers this direction while approaching the origin.

Concretely, define
\[
\widetilde{\AAA^2}
=
\left\{
\bigl((z_1,z_2),[s:t]\bigr)\in\AAA^2\times\PP^1
:
z_1t=z_2s
\right\}.
\]
The equation says that the vector \((z_1,z_2)\) lies on the line with
direction \([s:t]\). Projection onto the first factor gives a map
\[
\pi:\widetilde{\AAA^2}\longrightarrow\AAA^2.
\]

If \((z_1,z_2)\neq(0,0)\), there is a unique compatible direction:
\[
[s:t]=[z_1:z_2].
\]
Hence \(\pi\) is an isomorphism away from the origin. Over the origin,
however, the equation becomes
\[
0=0,
\]
so every direction is allowed:
\[
\pi^{-1}(0,0)\simeq\PP^1.
\]
The blowup therefore replaces the origin by a projective line
parameterizing directions through it. This projective line is called
the \emph{exceptional divisor} and will be denoted by \(E\).

\begin{figure}[ht]
\centering
\begin{tikzpicture}[scale=1.0]
  % Before.
  \begin{scope}[xshift=-4.3cm]
    \draw[->,black!65] (-2.1,0) -- (2.2,0);
    \draw[->,black!65] (0,-1.9) -- (0,2.1);
    \draw[thick,DictionaryBlue] (-1.8,-0.9) -- (1.8,0.9);
    \draw[thick,SurfaceGreen] (-1.0,-1.7) -- (1.0,1.7);
    \draw[thick,ExampleGold] (-1.8,1.15) -- (1.8,-1.15);
    \fill[WarningRed] (0,0) circle (2.5pt);
    \node[below] at (0,-2.15) {\(\AAA^2\)};
    \node[above right,WarningRed] at (0.1,0.1) {origin};
  \end{scope}

  \draw[very thick,-{Latex[length=2.8mm]},DictionaryBlue]
    (-1.25,0) -- (1.25,0)
    node[midway,above] {blow up};

  % After.
  \begin{scope}[xshift=4.3cm]
    \draw[->,black!65] (-2.1,0) -- (2.2,0);
    \draw[->,black!65] (0,-1.9) -- (0,2.1);

    \draw[very thick,magenta!70!DictionaryBlue]
      (0,-1.05) arc[start angle=-90,end angle=270,x radius=0.72,y radius=1.05];

    \draw[thick,DictionaryBlue] (-1.85,-0.92) -- (-0.52,-0.26);
    \draw[thick,DictionaryBlue] (0.52,0.26) -- (1.85,0.92);

    \draw[thick,SurfaceGreen] (-1.02,-1.72) -- (-0.28,-0.48);
    \draw[thick,SurfaceGreen] (0.28,0.48) -- (1.02,1.72);

    \draw[thick,ExampleGold] (-1.82,1.16) -- (-0.48,0.31);
    \draw[thick,ExampleGold] (0.48,-0.31) -- (1.82,-1.16);

    \node[right,magenta!70!DictionaryBlue] at (0.78,0.95)
      {\(E\simeq\PP^1\)};
    \node[below] at (0,-2.15) {\(\widetilde{\AAA^2}\)};
  \end{scope}
\end{tikzpicture}
\caption{The blowup replaces the origin by the projective line of
directions through it. The drawing is schematic: different directions
meet the exceptional divisor at different points.}
\label{fig:a2-blowup-directions}
\end{figure}

\begin{observation}
It is tempting to say that a blowup simply removes a point and inserts
a copy of \(\PP^1\). This is a useful picture, but the morphism
\(\pi\) carries more information: it explains how nearby points
approach the exceptional divisor and which limiting direction they
select.
\end{observation}

\subsection{The two affine charts of the blowup}

The projective line of directions has two standard charts, and these
give two affine charts on the blowup.

On the chart \(s\neq0\), write
\[
q=\frac{t}{s}.
\]
The equation \(z_1t=z_2s\) becomes
\[
z_2=z_1q.
\]
Thus this part of the blowup has coordinates \((z_1,q)\), and the map
to \(\AAA^2\) is
\[
(z_1,q)\longmapsto(z_1,z_1q).
\]
The exceptional divisor is given by
\[
z_1=0.
\]

On the chart \(t\neq0\), write
\[
p=\frac{s}{t}.
\]
The defining equation becomes
\[
z_1=pz_2.
\]
This chart has coordinates \((p,z_2)\), the blowup map is
\[
(p,z_2)\longmapsto(pz_2,z_2),
\]
and the exceptional divisor is given by
\[
z_2=0.
\]

\begin{figure}[ht]
\centering
\begin{tikzpicture}[
  chart/.style={
    rounded corners,
    minimum width=4.1cm,
    minimum height=2.0cm,
    align=center,
    draw=black!55
  }
]
  \node[chart,fill=PastelYellow] (A) at (-3.3,0)
    {\(\displaystyle s\neq0\)\\[1mm]
     \(\displaystyle z_2=z_1q\)\\
     coordinates \((z_1,q)\)\\
     \(E=\{z_1=0\}\)};

  \node[chart,fill=DictionaryBlue!8] (B) at (3.3,0)
    {\(\displaystyle t\neq0\)\\[1mm]
     \(\displaystyle z_1=pz_2\)\\
     coordinates \((p,z_2)\)\\
     \(E=\{z_2=0\}\)};

  \draw[thick,<->,magenta!70!DictionaryBlue]
    (A) -- node[above,fill=white,inner sep=2pt]
    {\(\displaystyle p=q^{-1}\)}
    (B);
\end{tikzpicture}
\caption{The two affine charts of the blowup. On their overlap, the
two coordinates \(p\) and \(q\) are reciprocal coordinates on the
exceptional projective line.}
\label{fig:a2-blowup-charts}
\end{figure}

The two coordinate axes through the origin are separated by the
blowup. In \(U_0\), the coordinate lines through \(p_0\) are
\[
D_1\cap U_0=\{z_1=0\},
\qquad
D_2\cap U_0=\{z_2=0\}.
\]
Their strict transforms no longer meet one another. Instead, each
meets \(E\) at the point corresponding to its own direction:
\[
\widetilde D_1\cap E=\{[s:t]=[0:1]\},
\qquad
\widetilde D_2\cap E=\{[s:t]=[1:0]\}.
\]

\begin{figure}[ht]
\centering
\begin{tikzpicture}[scale=1.0]
  % Before.
  \begin{scope}[xshift=-3.8cm]
    \draw[very thick,ExampleGold] (-2.0,0) -- (2.0,0);
    \draw[very thick,SurfaceGreen] (0,-1.7) -- (0,1.7);
    \fill[DictionaryBlue] (0,0) circle (2.4pt);
    \node[below] at (0,-2.05) {\(D_1\cap D_2=\{p_0\}\)};
  \end{scope}

  \draw[very thick,-{Latex[length=2.6mm]},DictionaryBlue]
    (-1.15,0) -- (1.15,0);

  % After.
  \begin{scope}[xshift=3.8cm]
    \draw[very thick,magenta!70!DictionaryBlue] (0,-1.15) -- (0,1.15);
    \draw[very thick,ExampleGold] (-2.0,0.78) -- (0,0.78);
    \draw[very thick,SurfaceGreen] (0,-0.78) -- (2.0,-0.78);

    \fill[DictionaryBlue] (0,0.78) circle (2.2pt);
    \fill[DictionaryBlue] (0,-0.78) circle (2.2pt);

    \node[left,ExampleGold!75!black] at (-1.2,0.78) {\(\widetilde D_1\)};
    \node[right,SurfaceGreen!70!black] at (1.2,-0.78) {\(\widetilde D_2\)};
    \node[right,magenta!70!DictionaryBlue] at (0.15,1.08) {\(E\)};
    \node[below] at (0,-2.05)
      {\(\widetilde D_1\cap\widetilde D_2=\varnothing\)};
  \end{scope}
\end{tikzpicture}
\caption{Before the blowup, \(D_1\) and \(D_2\) meet at \(p_0\).
Afterward, their strict transforms meet the exceptional divisor at
two different points.}
\label{fig:p2-blowup-separates-divisors}
\end{figure}

\subsection{The star subdivision of the fan}

The fixed point \(p_0\) corresponds to the cone
\[
\sigma_0=\Cone{v_1,v_2},
\qquad
v_1=(1,0),\quad v_2=(0,1).
\]
Insert the new primitive vector
\[
v_E=v_1+v_2=(1,1)
\]
and the ray
\[
\rho_E=\RR_{\geq0}v_E.
\]
The old cone \(\sigma_0\) is replaced by two cones:
\[
\Cone{v_1,v_E}
\qquad\text{and}\qquad
\Cone{v_E,v_2}.
\]
All other cones of the fan are left unchanged.

This operation is called the \emph{star subdivision} of
\(\sigma_0\) by the ray \(\rho_E\). The resulting fan describes the
blowup
\[
\pi:\Blp{\PP^2}{p_0}\longrightarrow\PP^2.
\]

\begin{figure}[ht]
\centering
\begin{tikzpicture}[scale=0.95]
  % Old fan.
  \begin{scope}[xshift=-4.1cm]
    \fill[PastelYellow] (0,0) -- (2.7,0) -- (0,2.7) -- cycle;
    \fill[DictionaryBlue!8] (0,0) -- (0,2.7) -- (-2.3,-2.3) -- cycle;
    \fill[SurfaceGreen!9] (0,0) -- (-2.3,-2.3) -- (2.7,0) -- cycle;

    \draw[very thick,WarningRed] (0,0) -- (-2.55,-2.55);
    \draw[very thick,ExampleGold] (0,0) -- (2.9,0);
    \draw[very thick,SurfaceGreen] (0,0) -- (0,2.9);

    \node at (1.05,1.0) {\(\sigma_0\)};
    \node[below] at (0,-3.0) {fan of \(\PP^2\)};
  \end{scope}

  \draw[very thick,-{Latex[length=2.8mm]},DictionaryBlue]
    (-1.2,0.2) -- (1.2,0.2)
    node[midway,above,align=center]
    {insert\\\(\rho_E\)};

  % New fan.
  \begin{scope}[xshift=4.1cm]
    \fill[PastelYellow] (0,0) -- (2.7,0) -- (1.9,1.9) -- cycle;
    \fill[magenta!12] (0,0) -- (1.9,1.9) -- (0,2.7) -- cycle;
    \fill[DictionaryBlue!8] (0,0) -- (0,2.7) -- (-2.3,-2.3) -- cycle;
    \fill[SurfaceGreen!9] (0,0) -- (-2.3,-2.3) -- (2.7,0) -- cycle;

    \draw[very thick,WarningRed] (0,0) -- (-2.55,-2.55);
    \draw[very thick,ExampleGold] (0,0) -- (2.9,0);
    \draw[very thick,SurfaceGreen] (0,0) -- (0,2.9);
    \draw[very thick,magenta!70!DictionaryBlue] (0,0) -- (2.25,2.25);

    \node[below] at (0,-3.0)
      {fan of \(\Blp{\PP^2}{p_0}\)};
    \node[right,magenta!70!DictionaryBlue] at (1.72,1.72)
      {\(\rho_E\)};
  \end{scope}
\end{tikzpicture}
\caption{The star subdivision replaces the cone for \(p_0\) by two
cones separated by the new ray \(\rho_E\).}
\label{fig:p2-star-subdivision}
\end{figure}

Why should the new vector be \(v_1+v_2\)? In the local chart
\(U_0\simeq\AAA^2\), the rays generated by \(v_1\) and \(v_2\)
represent the two coordinate directions. The vector \((1,1)\)
records approaching the origin while both coordinates vanish at the
same rate:
\[
t\longmapsto(t,t).
\]
It is the simplest integral direction lying strictly between the two
old rays.

\subsection{The exceptional divisor in the orbit picture}

The new ray \(\rho_E\) corresponds to a new one-dimensional torus
orbit and its closure, the exceptional divisor \(E\). The two new
maximal cones correspond to the two torus-fixed points on \(E\):
\[
\Cone{v_1,v_E}
\longleftrightarrow
\widetilde D_1\cap E,
\]
and
\[
\Cone{v_E,v_2}
\longleftrightarrow
E\cap\widetilde D_2.
\]

The old fixed point \(p_0\) has disappeared. In its place we find the
one-dimensional orbit inside \(E\) together with two new fixed points.
The orbit count changes from
\[
1+3+3=7
\]
to
\[
1+4+4=9:
\]
one open orbit, four one-dimensional orbits, and four fixed points.

\begin{figure}[ht]
\centering
\begin{tikzpicture}[scale=1.0]
  \coordinate (A) at (-2.7,0);
  \coordinate (B) at (2.7,0);
  \coordinate (C) at (1.15,2.7);
  \coordinate (D) at (-1.15,2.7);

  \fill[PastelYellow] (-1.05,2.15) -- (-1.95,0.42) --
    (1.95,0.42) -- (1.05,2.15) -- cycle;

  \draw[very thick,WarningRed] (A) -- (B);
  \draw[very thick,SurfaceGreen] (B) -- (C);
  \draw[very thick,magenta!70!DictionaryBlue] (C) -- (D);
  \draw[very thick,ExampleGold] (D) -- (A);

  \fill[DictionaryBlue] (A) circle (2.3pt);
  \fill[DictionaryBlue] (B) circle (2.3pt);
  \fill[DictionaryBlue] (C) circle (2.3pt);
  \fill[DictionaryBlue] (D) circle (2.3pt);

  \node at (0,1.15) {\(T\)};
  \node[above,magenta!70!DictionaryBlue] at (0,2.75) {\(E\)};
  \node[below] at (0,-0.45)
    {boundary incidence after blowing up \(p_0\)};
\end{tikzpicture}
\caption{The boundary incidence pattern after the blowup is a
quadrilateral. The new side represents the exceptional divisor, and
its two endpoints represent the two new fixed points.}
\label{fig:p2-blowup-orbit-quadrilateral}
\end{figure}

\begin{observation}
The surface \(\Blp{\PP^2}{p_0}\) is isomorphic to the first Hirzebruch
surface \(\Hirz{1}\). Hirzebruch surfaces form the next natural family
of smooth complete toric surfaces after \(\PP^2\), and they will give
us a useful supply of examples.
\end{observation}

\subsection{The exceptional curve has self-intersection \texorpdfstring{\(-1\)}{-1}}

One numerical feature of the exceptional divisor is worth previewing:
\[
E^2=-1.
\]
Intersection theory comes later, so this formula is a preview rather
than a completed calculation.

Geometrically, the negative self-intersection records that the
exceptional curve can be contracted back to a smooth point. The
blowup map
\[
\pi:\Blp{\PP^2}{p_0}\longrightarrow\PP^2
\]
does exactly this: it sends the whole curve \(E\) to \(p_0\) and is an
isomorphism elsewhere. In the fan, contraction means deleting the ray
\(\rho_E\) and restoring the original cone
\(\Cone{v_1,v_2}\).

\begin{threedictionaries}
The blowup of \(p_0\) is one operation described in three languages.

\medskip

\noindent
\begin{minipage}[t]{0.30\linewidth}
\centering
\textbf{Algebra}\par
\smallskip
Locally, blow up the ideal
\[
(z_1,z_2)
\]
at the origin of
\(\Spec\Bbbk[z_1,z_2]\).
The two charts satisfy
\[
z_2=z_1q,\qquad z_1=pz_2.
\]
\end{minipage}
\hfill
\begin{minipage}[t]{0.30\linewidth}
\centering
\textbf{Geometry}\par
\smallskip
Replace \(p_0\) by the exceptional curve
\[
E\simeq\PP^1,
\]
whose points record directions through \(p_0\).
\end{minipage}
\hfill
\begin{minipage}[t]{0.30\linewidth}
\centering
\textbf{Combinatorics}\par
\smallskip
Insert the ray
\[
\rho_E=\RR_{\geq0}(v_1+v_2)
\]
and subdivide
\(\Cone{v_1,v_2}\) into two cones.
\end{minipage}
\end{threedictionaries}

\begin{surfacelaboratory}
Track the blowup of \(p_0\) through the following parallel changes.

\begin{center}
\small
\begin{tabularx}{\linewidth}{@{}
  >{\raggedright\arraybackslash}X
  >{\raggedright\arraybackslash}X
  >{\raggedright\arraybackslash}X@{}}
\toprule
\textbf{Before}
&
\textbf{Blowup operation}
&
\textbf{After}
\\
\midrule
the ideal \((z_1,z_2)\)
&
record the ratio \([z_1:z_2]\)
&
two affine blowup charts
\\[1mm]
the point \(p_0=D_1\cap D_2\)
&
replace the point by its directions
&
\(E\simeq\PP^1\), meeting
\(\widetilde D_1,\widetilde D_2\)
\\[1mm]
the cone \(\Cone{v_1,v_2}\)
&
insert \(v_E=v_1+v_2\)
&
two cones and one new ray.
\\
\bottomrule
\end{tabularx}
\end{center}
\end{surfacelaboratory}


\section*{What this chapter has accomplished}
\addcontentsline{toc}{section}{What this chapter has accomplished}
\markboth{What this chapter has accomplished}{What this chapter has accomplished}
\label{sec:p2-predictions}

We have spent an entire chapter revisiting a single familiar surface.
The point was not to make \(\PP^2\) look more complicated. It was to
notice that several kinds of information we usually study separately
are shadows of the same combinatorial structure.

In \(\PP^2\), we have been able to check everything directly. The
next chapters turn these observations into definitions and theorems
for arbitrary toric varieties.

\subsection*{The affine dictionary}

For the affine plane, the first quadrant
\[
\sigma=\RR_{\geq0}^2\subseteq N_{\RR}
\]
controlled two complementary questions:

\begin{itemize}
\item which monomial curves acquire limits in \(\AAA^2\);
\item which Laurent monomials extend to regular functions on
\(\AAA^2\).
\end{itemize}

The corresponding algebra was
\[
\Bbbk[\sigma^\vee\cap M]
=
\Bbbk[x,y],
\]
and its spectrum was the affine plane.

This suggests the general construction
\[
\sigma
\quad\rightsquigarrow\quad
\sigma^\vee\cap M
\quad\rightsquigarrow\quad
\Bbbk[\sigma^\vee\cap M]
\quad\rightsquigarrow\quad
\Spec\Bbbk[\sigma^\vee\cap M].
\]
The general construction developed in the next chapters will show
that rational polyhedral cones produce affine toric varieties in
precisely this way.

\subsection*{Fans as gluing instructions}

The fan of \(\PP^2\) consisted of three maximal cones. Each maximal
cone produced one affine chart, and shared faces produced the overlap
data. This suggests that a fan should be understood as a set of
instructions for gluing affine toric varieties.

The orbit decomposition was visible in the same fan:
\[
\{0\}\longleftrightarrow T,\qquad
\rho_i\longleftrightarrow O_i,\qquad
\sigma_i\longleftrightarrow p_i.
\]
The dimensions added to \(2\):
\[
\dim\sigma+\dim O(\sigma)=2.
\]
In higher dimensions, the same pattern will become the
orbit--cone correspondence.

\subsection*{Polytopes, sections, and projective embeddings}

The dilated triangle \(d\Delta\) indexed the monomial sections of
\(\cO_{\PP^2}(d)\):
\[
x_0^{d-a-b}x_1^ax_2^b
\longleftrightarrow
(a,b)\in d\Delta\cap M.
\]
Counting global sections became counting lattice points.

This suggests that lattice polytopes should encode projective
geometry. Their lattice points should give monomial maps to projective
space, their faces should record boundary behavior, and their normal
fans should construct projective toric varieties.

\begin{warningbox}
Fans are more general than normal fans of polytopes. Every lattice
polytope has a normal fan, but not every fan comes from a polytope.
Consequently, not every toric variety is projective.
\end{warningbox}

\subsection*{Rays, divisors, and birational modifications}

In the fan of \(\PP^2\), each ray corresponded to one coordinate line:
\[
\rho_i\longleftrightarrow D_i.
\]
After blowing up \(p_0\), the new ray corresponded to the exceptional
divisor:
\[
\rho_E\longleftrightarrow E.
\]
This predicts a general correspondence between rays of a fan and
torus-invariant prime divisors.

The star subdivision also suggests that combinatorial refinements
should produce geometric modifications:
\[
\text{subdivide a fan}
\quad\rightsquigarrow\quad
\text{modify a toric variety birationally}.
\]
For a smooth toric surface, inserting the ray generated by the sum of
two adjacent primitive ray generators describes the blowup of the
corresponding fixed point. When the reverse operation is allowed,
deleting that ray contracts the exceptional curve.

\subsection*{What we have seen, and what remains to prove}

The following table separates the evidence supplied by \(\PP^2\) from
the general theory still ahead.

\begin{center}
\small
\begin{tabularx}{\linewidth}{@{}
  >{\raggedright\arraybackslash}X
  >{\raggedright\arraybackslash}X@{}}
\toprule
\textbf{Seen directly in \(\PP^2\)}
&
\textbf{General principle to be developed}
\\
\midrule
The first-quadrant cone gives \(\AAA^2\).
&
A rational polyhedral cone gives an affine toric variety.
\\[1mm]
The three maximal cones give the three affine charts.
&
A fan glues affine toric varieties into a global toric variety.
\\[1mm]
Cones index the seven torus orbits.
&
Cones of a fan correspond to torus orbits, with complementary
dimensions.
\\[1mm]
The three rays give the coordinate lines.
&
Rays correspond to torus-invariant prime divisors.
\\[1mm]
Lattice points in \(d\Delta\) index sections of
\(\cO_{\PP^2}(d)\).
&
Lattice polytopes and polyhedra encode sections of toric line bundles.
\\[1mm]
The normal fan of \(\Delta\) recovers \(\PP^2\).
&
Normal fans connect lattice polytopes with projective toric varieties.
\\[1mm]
Adding the ray \((1,1)\) blows up \(p_0\).
&
Under suitable hypotheses, star subdivisions describe toric blowups
of invariant strata.
\\[1mm]
Deleting the exceptional ray contracts \(E\).
&
Fan refinements and contractions organize toric birational geometry.
\\
\bottomrule
\end{tabularx}
\end{center}

The chapter's organizing principle may now be displayed in one
picture. Algebra, geometry, and combinatorics are not three separate
topics: they are three ways of reading the same constructions.

\begin{figure}[ht]
\centering
\begin{tikzpicture}[
  mainbox/.style={
    rounded corners,
    text width=3.0cm,
    minimum height=1.65cm,
    align=center,
    draw=black!55,
    font=\small
  }
]
  \node[mainbox,fill=DictionaryBlue!8] (alg) at (-4.0,0.25)
    {\textbf{Algebra}\\[1mm]
     monomials and semigroup rings\\
     gradings and ideals};

  \node[mainbox,fill=PastelYellow] (geo) at (0,0.25)
    {\textbf{Geometry}\\[1mm]
     charts, orbits, and divisors\\
     line bundles and blowups};

  \node[mainbox,fill=SurfaceGreen!9] (comb) at (4.0,0.25)
    {\textbf{Combinatorics}\\[1mm]
     lattices, cones, and fans\\
     polytopes and subdivisions};

  \draw[very thick,<->,ExampleGold]
    (alg.east) -- (geo.west);
  \draw[very thick,<->,magenta!70!DictionaryBlue]
    (geo.east) -- (comb.west);

  \draw[very thick,<->,WarningRed]
    (alg.south) to[out=-35,in=-145]
    node[midway,below=4mm,fill=white,inner sep=2pt]
    {the three dictionaries}
    (comb.south);
\end{tikzpicture}
\caption{The organizing principle of the course. None of the three
languages is merely decorative: each makes some constructions
transparent that are difficult to see in the other two.}
\label{fig:p2-course-dictionaries}
\end{figure}

\begin{threedictionaries}
The chapter may be compressed into four recurring translations.

\medskip

\begin{center}
\small
\begin{tabularx}{\linewidth}{@{}
  >{\raggedright\arraybackslash}X
  >{\centering\arraybackslash}p{1.0cm}
  >{\raggedright\arraybackslash}X
  >{\centering\arraybackslash}p{1.0cm}
  >{\raggedright\arraybackslash}X@{}}
monomial functions
&
\(\longleftrightarrow\)
&
regular functions on torus charts
&
\(\longleftrightarrow\)
&
lattice points in dual cones
\\[2mm]
coordinate rings
&
\(\longleftrightarrow\)
&
affine toric charts
&
\(\longleftrightarrow\)
&
cones of a fan
\\[2mm]
graded monomials
&
\(\longleftrightarrow\)
&
sections of line bundles
&
\(\longleftrightarrow\)
&
lattice points in polytopes
\\[2mm]
blowup ideals
&
\(\longleftrightarrow\)
&
exceptional divisors
&
\(\longleftrightarrow\)
&
star subdivisions.
\end{tabularx}
\end{center}
\end{threedictionaries}

\subsection*{A provisional picture of a toric variety}

The examples suggest the following working picture. A toric variety
is an algebraic variety containing an algebraic torus as a dense open
subset, together with an action of that torus extending its
multiplication on itself.

This is not yet our formal starting point. We still need to understand
algebraic tori, their monomial functions, and their one-parameter
motions in a coordinate-free way. That is the task of the next
chapter.

\begin{guidingquestion}
Which features of \(\Gm^n\) are intrinsic, and which arose only
because we chose coordinates \(x_1,\ldots,x_n\)?
\end{guidingquestion}


\clearpage

\begin{exercises}[Comprehension exercises]
\markboth{Comprehension exercises}{Comprehension exercises}

\item Determine which characters of \(T=\Gm^2\) extend to regular
functions on \(\AAA^2\).

\item Let \(\oneps{(u,v)}(t)=(t^u,t^v)\). Determine when this map
extends to a morphism \(\AAA^1\to\AAA^2\).

\item Write the transition maps between the three standard affine
charts of \(\PP^2\).

\item List all torus orbits in \(\PP^2\), compute their dimensions, and
describe their closures.

\item For \(d=2\) and \(d=3\), draw \(d\Delta\cap M\) and label each
lattice point by its corresponding monomial section of
\(\cO_{\PP^2}(d)\).

\item Compute
\[
\dim_{\Bbbk}\Hzero{\PP^2}{\cO_{\PP^2}(d)}
\]
by counting lattice points in \(d\Delta\).

\item Compute the inward-pointing primitive normal vector to each edge
of \(\Delta\), and use these vectors to draw the normal fan.

\item Draw the star subdivision corresponding to the blowup of each
of the three torus-fixed points of \(\PP^2\).

\item Explain why blowing up any one of the three torus-fixed points
produces isomorphic toric surfaces.

\end{exercises}

\clearpage
\section*{Challenge exercise}
\addcontentsline{toc}{section}{Challenge exercise}
\markboth{Challenge exercise}{Challenge exercise}

\noindent
\textbf{Compute the topological Euler characteristic of
\(\PP^2(\CC)\) using the torus action.}

\medskip

The goal is to discover a map from \(\PP^2(\CC)\) to the triangle
\[
\Delta=\Conv\{(0,0),(1,0),(0,1)\}
\]
whose fibers reflect the orbit decomposition. Do not use the usual
cell decomposition
\(\PP^2=\CC^2\sqcup\CC\sqcup\{\mathrm{pt}\}\).

\begin{enumerate}[label=\textbf{Hint \arabic*.},wide=0pt]

\item Consider the map
\[
\mu:\PP^2(\CC)\longrightarrow\RR^2,
\qquad
\mu([z_0:z_1:z_2])
=
\left(
\frac{|z_1|^2}{|z_0|^2+|z_1|^2+|z_2|^2},
\frac{|z_2|^2}{|z_0|^2+|z_1|^2+|z_2|^2}
\right).
\]
Check that it is well defined on projective points and that its image
is exactly \(\Delta\). This is the moment map for the compact torus
\((S^1)^2\subset(\CC^\times)^2\), up to the usual harmless choices of
normalization.

\item Let \((a,b)\) lie in the interior of \(\Delta\). Fixing
\(\mu([z_0:z_1:z_2])=(a,b)\) fixes the absolute values of the three
homogeneous coordinates, up to common scaling. Determine which phase
parameters remain. Show that the fiber is a real two-torus:
\[
\mu^{-1}(a,b)\simeq(S^1)^2.
\]

\item Now let \((a,b)\) lie in the relative interior of an edge of
\(\Delta\). One homogeneous coordinate must vanish. Show that the
fiber is a circle:
\[
\mu^{-1}(a,b)\simeq S^1.
\]

\item Determine the fibers over the three vertices of \(\Delta\).
Show that they are the three torus-fixed points
\[
p_0,\qquad p_1,\qquad p_2.
\]

\item Stratify \(\Delta\) into the relative interiors of its face, its
three edges, and its three vertices. Over each stratum, \(\mu\) is a
locally trivial fibration of the type found above. Use additivity of
Euler characteristic over a finite stratification and multiplicativity
for locally trivial fibrations. On the noncompact strata, one may use
compactly supported Euler characteristic; since \(\PP^2(\CC)\) is
compact, the final result agrees with its ordinary Euler
characteristic.

\item Recall that
\[
\chi\bigl((S^1)^2\bigr)=0,
\qquad
\chi(S^1)=0,
\qquad
\chi(\{\mathrm{pt}\})=1.
\]
Conclude that only the fibers above the three vertices contribute.
What is
\[
\chi\bigl(\PP^2(\CC)\bigr)?
\]

\item Now forget the special geometry of \(\PP^2\) for a moment and
look only at its decomposition into torus orbits. Every
positive-dimensional orbit is a copy of
\((\CC^\times)^r\) for some \(r>0\), and hence has Euler
characteristic zero. What conjecture does this suggest for the Euler
characteristic of an arbitrary complex toric variety? For a
noncomplete variety, you may formulate the conjecture first using
compactly supported Euler characteristic. Can you express your answer
in terms of torus-fixed points, or equivalently in terms of the
full-dimensional cones of the fan?

\end{enumerate}


\section*{Notes and references}
\addcontentsline{toc}{section}{Notes and references}
\markboth{Notes and references}{Notes and references}

This chapter is intentionally panoramic. It uses \(\PP^2\) to preview
the constructions that will be defined and proved in the chapters
that follow: affine toric varieties from cones, global varieties from
fans, invariant divisors from rays, monomial sections from lattice
points, and toric blowups from star subdivisions.

For systematic introductions to these correspondences, see Fulton,
\emph{Introduction to Toric Varieties}, Chapters~1--3, and
Cox--Little--Schenck, \emph{Toric Varieties}, Chapters~1--4 and~6.
The moment-map viewpoint used in the challenge is developed from
first principles in Chapter~8 of these notes.
